\documentclass{article}
\usepackage{amssymb}
\usepackage{amsmath}
\usepackage{amsthm}
\usepackage{tikz-cd}
\usepackage{mathpartir}
\usepackage{stmaryrd}
\usepackage{parskip}
\usepackage{xcolor}
\usepackage{array}
\usepackage[margin=1.0in]{geometry}
\usepackage{comment}
\usepackage{quiver}

\definecolor{myorange}{RGB}{242, 114, 0}
\definecolor{myblue}{RGB}{0,162,255}



%%%%%%%%%%%%%%%%%%%%%%%%
% theorem environments
%%%%%%%%%%%%%%%%%%%%%%%%
\newtheorem{theorem}{Theorem}[section]
\newtheorem{lemma}{Lemma}[section]
\newtheorem{nonnum-theorem}{Theorem}[section]
\newtheorem{corollary}{Corollary}[section]

\theoremstyle{definition}
\newtheorem{definition}{Definition}[section]

\theoremstyle{remark}
\newtheorem{remark}{Remark}

\newtheorem{example}{Example}
%%%%%%%%%%%%%%%%%%%%%%%%%%%%%%%%%%%%%%%%%%%%%%%%%%%

\newcommand{\To}{\Rightarrow}
\newcommand{\inl}{\mathsf{inl}}
\newcommand{\inr}{\mathsf{inr}}
\newcommand{\alt}{\mathrel{\bf \,\mid\,}}

\newcommand{\ob}{\text{Ob}}


\newcommand{\extlc}{\text{Ext-}\lambda}
\newcommand{\extlcm}{\text{Ext-}\lambda^{-\text{trans}}}
\newcommand{\extlcmm}{\text{Ext-}\lambda^{-\text{trans}-\text{cast}}}
\newcommand{\extlcprime}{\text{Ext-}\lambda'}
\newcommand{\intlc}{\text{Int-}\lambda}
\newcommand{\intlcbisim}{\text{Int$_\approx$-}\lambda}
\newcommand{\erase}[1]{\lfloor {#1} \rfloor}


\newcommand{\uarrowl}{\mathrel{\rotatebox[origin=c]{-30}{$\leftarrowtail$}}}
\newcommand{\uarrowr}{\mathrel{\rotatebox[origin=c]{60}{$\leftarrowtail$}}}
\newcommand{\darrowl}{\mathrel{\rotatebox[origin=c]{30}{$\twoheadleftarrow$}}}
\newcommand{\darrowr}{\mathrel{\rotatebox[origin=c]{120}{$\twoheadleftarrow$}}}
\newcommand{\vuarrow}{\mathrel{\rotatebox[origin=c]{-90}{$\leftarrowtail$}}}
\newcommand{\vdarrow}{\mathrel{\rotatebox[origin=c]{90}{$\twoheadleftarrow$}}}

% Types, terms, and precision 

\newcommand{\dyn}{?}
\newcommand{\nat}{\text{Nat}}
\newcommand{\bool}{\text{Bool}}
\newcommand{\ra}{\rightharpoonup}
\newcommand{\Ret}[1]{\mathsf{Ret\,}{#1}}
\newcommand{\hole}[1]{\bullet \colon {#1}}
\newcommand{\dyntodyn}{\dyn \ra\, \dyn}


\newcommand{\up}[2]{\langle{#2}\uarrowl{#1}\rangle}
\newcommand{\dn}[2]{\langle{#1}\darrowl{#2}\rangle}

\newcommand{\upc}[1]{\text{up}\,{#1}\,}
\newcommand{\dnc}[1]{\text{dn}\,{#1}\,}


% \newcommand{\ret}{\mathsf{ret}}
\newcommand{\err}{\mho}
\newcommand{\zro}{\textsf{zro}}
\newcommand{\suc}{\textsf{suc}}
\newcommand{\lda}[2]{\lambda {#1} . {#2}}

\newcommand{\injarr}[1]{\textsf{Inj}_\ra ({#1})}
\newcommand{\injnat}[1]{\textsf{Inj}_\text{nat} ({#1})}
\newcommand{\casenat}[4]{\text{Case}_\text{nat} ({#1}) \{ \text{no} \to {#2} \alt \text{nat}({#3}) \to {#4} \}}
\newcommand{\casearr}[4]{\text{Case}_\ra ({#1}) \{ \text{no} \to {#2} \alt \text{fun}({#3}) \to {#4} \}}
\newcommand{\casedyn}[5]{\text{Case}_\Dyn ({#1}) \{ \text{nat}({#2}) \to {#3} \alt \text{fun}({#4}) \to {#5} \}}
%\newcommand{\bind}[3]{\text{var } {#1} = {#2} \text{ in } {#3}}
\newcommand{\matchnat}[4]{\text{match } {#1} \text{ with } \{ \textsf {zero} \Rightarrow {#2} \alt \textsf{ suc } {#3} \Rightarrow {#4} \}}

\newcommand{\refl}{\text{refl}}

\newcommand{\Lift}{\text{Lift}}

%\newcommand{\rel}{\circ\hspace{-4px}-\hspace{-4px}\bullet}
\newcommand{\rel}{\mathrel{\circ\mkern-6mu-\mkern-6mu\bullet}}
% \newcommand{\wand}{\mathrel{-\mkern-6mu*}}
\newcommand{\wand}{\mathrel{\multimap}}
\newcommand{\ltdyn}{\sqsubseteq}
\newcommand{\gtdyn}{\sqsupseteq}
\newcommand{\equidyn}{\mathrel{\gtdyn\ltdyn}}
\newcommand{\gamlt}{\Gamma^\ltdyn}
\newcommand{\deltalt}{\Delta^\ltdyn}
\newcommand{\relcomp}{\odot}

\newcommand{\hasty}[3]{{#1} \vdash {#2} \colon {#3}}
\newcommand{\vhasty}[3]{{#1} \vdash^v {#2} \colon {#3}}
\newcommand{\phasty}[3]{{#1} \vdash^c {#2} \colon {#3}}
\newcommand{\etmprec}[4]{{#1} \vdash {#2} \ltdyn_e {#3} \colon {#4}}
\newcommand{\itmprec}[4]{{#1} \vdash {#2} \ltdyn_i {#3} \colon {#4}}
\newcommand{\etmequidyn}[4]{{#1} \vdash {#2} \equidyn_e {#3} \colon {#4}}
\newcommand{\itmequidyn}[4]{{#1} \vdash {#2} \equidyn_i {#3} \colon {#4}}
\newcommand{\synbisim}{\approx_{\text{syn}}}

\newcommand{\Dwn}{\Downarrow}
\newcommand{\qte}[1]{\text{quote}({#1})}

\newcommand{\elab}[1]{\text{Elab}({#1})}

% Perturbations
\newcommand{\pertp}{\text{Pert}^\text{P}}
\newcommand{\perte}{\text{Pert}^\text{E}}
\newcommand{\pertdyn}[2]{\text{pert-dyn}({#1}, {#2})}
\newcommand{\delaypert}[1]{\text{delay-pert}({#1})}

\newcommand{\pertc}{\text{Pert}_{\text{C}}}
\newcommand{\pertv}{\text{Pert}_{\text{V}}}


% SGDT and Intensional Stuff

\newcommand{\later}{{\vartriangleright}}
\newcommand{\laterhs}{{\later}}
\newcommand{\type}{\texttt{Type}}
\newcommand{\lob}{\text{L\"{o}b}}
\newcommand{\tick}{\mathsf{tick}}
\newcommand{\nxt}{\mathsf{next}}
\newcommand{\fix}{\mathsf{fix}}
\newcommand{\kpa}{\kappa}

% Model-related stuff
\newcommand{\calV}{\mathcal{V}}
\newcommand{\calE}{\mathcal{E}}
\newcommand{\calVk}{\mathcal{V}_k}
\newcommand{\calEk}{\mathcal{E}_k}
\newcommand{\tok}{\mathrel{\mathop{\to}\limits^{\textrm{\footnotesize k}}}}
\newcommand{\timesk}{\mathrel{\mathop{\times}\limits^{\textrm{k}}}}
\newcommand{\Fk}{\mathrel{\mathop{F}\limits^{\textrm{k}}}}
\newcommand{\Uk}{\mathrel{\mathop{U}\limits^{\textrm{k}}}}

\newcommand{\op}[1]{{#1}^{\textrm{op}}}
\newcommand{\calC}{\mathcal{C}}
\newcommand{\Set}{\mathsf{Set}}
\newcommand{\ErrDom}{\mathsf{ErrDom}}
\newcommand{\Yo}{\mathsf{Yo}}
\newcommand{\Hom}{\mathsf{Hom}}
\newcommand{\calS}{\mathcal{S}}
\newcommand{\gfix}{\texttt{gfix}}
\newcommand{\qfix}{\texttt{qfix}}
\newcommand{\calU}{\mathcal{U}}
\newcommand{\laterhat}{\widehat{\later}}
\newcommand{\El}{\mathsf{El}}
\newcommand{\Clock}{\mathsf{Clock}}

\newcommand{\Machine}[1]{\mathsf{Machine}\, {#1}}


% Predomains and EP pairs
\newcommand{\Nat}{\mathsf{Nat}}
\newcommand{\Dyn}{\mathsf{Dyn}}
\newcommand{\ty}[1]{\langle {#1} \rangle}
\newcommand{\li}{L_\mho}
\newcommand{\liclk}[1]{L_\mho [{#1}]}

\newcommand{\ext}[2]{\text{ext}\,{#1}\,{#2}}
\newcommand{\map}[2]{\text{map}\,{#1}\,{#2}}

\newcommand{\ltls}{\ltdyn}
\newcommand{\bisim}{\approx}
\newcommand{\semlt}{\le}
\newcommand{\semltbad}{\lesssim}

%\newcommand{\injarr}{\textsf{Inj}_\to}
%\newcommand{\injnat}{\textsf{Inj}_\mathbb{N}}

\newcommand{\id}{\mathsf{id}}
\newcommand{\ep}{\leadsto}

\newcommand{\emb}[2]{\mathsf{emb}_{#1}({#2})}
\newcommand{\proj}[2]{\mathsf{proj}_{#1}({#2})}

\newcommand{\monto}{\to_m}


% Notation for wait functions
\newcommand{\wre}{w_r^e}
\newcommand{\wle}{w_l^e}
\newcommand{\wrp}{w_r^p}
\newcommand{\wlp}{w_l^p}

\newcommand{\sem}[1]{\llbracket {#1} \rrbracket}
\newcommand{\semgl}[1]{{\sem{#1}}^\text{gl}}


% Denotational model
\newcommand{\errdom}{\textsf{ErrDom}}

\newcommand{\ptb}{\text{ptb}}
\newcommand{\push}{\text{push}}
\newcommand{\pull}{\text{pull}}

\newcommand{\pv}{P^{\mathcal{V}}}
\newcommand{\pe}{P^{\mathcal{E}}}
\newcommand{\ptbv}{\text{ptb}^\mathcal{V}}
\newcommand{\ptbe}{\text{ptb}^\mathcal{E}}

\newcommand{\Ppure}{P}
\newcommand{\Pk}{P^K}
\newcommand{\ptbk}{\text{ptb}^K}


\newcommand{\upf}{\text{up}}
\newcommand{\dnf}{\text{dn}
}
\newcommand{\arr}{\to}
\newcommand{\comp}{\bullet}

\newcommand{\ltsq}[2]{\mathrel{\ltdyn^{#1}_{#2}}}
\newcommand{\ltsqbisim}[2]{\mathrel{{\widetilde{\ltdyn}}^{#1}_{#2}}}
\newcommand{\ltbisim}{\mathrel{\widetilde{\ltdyn}}}
\newcommand{\vf}{\mathcal{V}_f}
\newcommand{\vsq}{\mathcal{V}_{sq}}
\newcommand{\ef}{\mathcal{E}_f}
\newcommand{\esq}{\mathcal{E}_{sq}}

\newcommand{\morbisimid}[1]{\text{Endo}_\bisim{#1}}


\newcommand{\sv}{s_{\mathcal{V}}}
\newcommand{\tv}{t_{\mathcal{V}}}
\newcommand{\rv}{r_{\mathcal{V}}}
\newcommand{\se}{s_{\mathcal{E}}}
\newcommand{\te}{t_{\mathcal{E}}}
\newcommand{\re}{r_{\mathcal{E}}}

\newcommand{\Ff}{F_f}
\newcommand{\Fsq}{F_{sq}}
\newcommand{\Uf}{U_f}
\newcommand{\Usq}{U_{sq}}
\newcommand{\arrf}{\arr_f}
\newcommand{\arrsq}{\arr_{sq}}

\newcommand{\vsim}{\mathcal{V}_{\bisim}}
\newcommand{\esim}{\mathcal{E}_{\bisim}}
\newcommand{\svsim}{s_{\mathcal{V}}^\bisim}
\newcommand{\tvsim}{t_{\mathcal{V}}^\bisim}
\newcommand{\rvsim}{r_{\mathcal{V}}^\bisim}
\newcommand{\sesim}{s_{\mathcal{E}}^\bisim}
\newcommand{\tesim}{t_{\mathcal{E}}^\bisim}
\newcommand{\resim}{r_{\mathcal{E}}^\bisim}



\newcommand{\ve}{\mathcal{V}_e}
\newcommand{\ee}{\mathcal{E}_e}

\newcommand{\vr}{\mathcal{V}_r}
\newcommand{\er}{\mathcal{E}_r}

\newcommand{\relatedin}[3]{{#1}\, {#3}\, {#2}}
\newcommand{\binrel}[1]{\mathbin{#1}}


\newcommand{\da}{\downarrow}

\newcommand{\ret}{\text{ret}}
\newcommand{\bind}[3]{{#1} <- {#2}\, ; \, {#3}}

\newcommand{\piv}{\Pi^\mathcal{V}}
\newcommand{\pie}{\Pi^\mathcal{E}}

\newcommand{\upl}{\textsc{UpL}}
\newcommand{\upr}{\textsc{UpR}}
\newcommand{\dnl}{\textsc{DnL}}
\newcommand{\dnr}{\textsc{DnR}}

\newcommand{\delre}{\delta^{r,e}}
\newcommand{\delle}{\delta^{l,e}}
\newcommand{\delrp}{\delta^{r,p}}
\newcommand{\dellp}{\delta^{l,p}}

\newcommand{\qordeq}{\bisim}

\newcommand{\inat}{\text{Inj}_\mathbb{N}}
\newcommand{\itimes}{\text{Inj}_\times}
\newcommand{\iarr}{\text{Inj}_\to}

\newcommand{\tnat}{\mathsf{nat}}
\newcommand{\ttimes}{\mathsf{times}}
\newcommand{\tfun}{\mathsf{fun}}

\newcommand{\cased}[7]{
    \begin{align*}
    \text{Case}_D ({#1}) &\text{ of } \{ \\
        &\alt \text{nat}({#2}) \to {#3}  \\
        &\alt \text{times}({#4}) \to {#5}  \\
        &\alt \text{fun}({#6}) \to {#7} \}
    \end{align*}
}

\newcommand{\inone}{\mathsf{in}_1}
\newcommand{\intwo}{\mathsf{in}_2}
\newcommand{\inthree}{\mathsf{in}_3}
\newcommand{\infour}{\mathsf{in}_4}
\newcommand{\infive}{\mathsf{in}_5}


\newcommand{\delay}{\text{Delay}}
\newcommand{\ledelay}{\le^{\text{Del}}}
\newcommand{\bisimdelay}{\bisim^{\text{Del}}}
\newcommand{\tnow}{\mathsf{now}}
\newcommand{\tlater}{\mathsf{later}}
\newcommand{\Da}{\Downarrow}




% Is there a notion of naturality for functors on double categories that
% involves horizontal morphisms?

% Come up with consistent notation for the double category of values of a double
% CBPV model, and likewise for computations




\begin{document}

\title{An Abstract, Guarded Categorical Model of Gradual Typing}
\author{Eric Giovannini}

\maketitle

\section{Introduction}

The work of Giovannini et al. \cite{gtt-sgdt} constructs a denotational
semantics for a simple gradually-typed lambda calculus using guarded type
theory.
%
Many of the aspects of the model construction are in fact independent of the
details of the particular language being modeled, e.g., the specific types and
effects in the language.
%
In this work, we seek to abstract over these details and present an abstract
categorical framework for constructing models of gradual typing in guarded type
theory.

With this abstract framework in hand, we then define concrete semantics for
gradually-typed languages with simple effects.

\subsection{Why an Abstract Model?}

The model developed by Giovannini et al. suffers from the limitation that it is
a single model for a quite bare-bones gradually-typed language. Real-world
languages can include effects, polymorphism, dependent types, etc. The goals of
an abstract model are twofold:

\begin{enumerate}
  \item To promote a better understanding of the essence of any guarded
  relational semantics of gradual tying independent of the particular language
  features being modeled.
  \item To allow for developers of languages to establish graduality for their
  language by instantiating our abstract framework.
\end{enumerate}

\subsection{Challenges}

Prior work by New and Licata \cite{new-licata18} defines a categorical model for
gradual typing that serves as the basis for denotational semantics using
\emph{classical} domain theory. Their model is based on \emph{double categories}
\cite{shulman2007}.
%
Unfortunately, as Giovannini et al. point out, adapting this model to the
guarded setting poses significant challenges, and much of their work consists of
defining a version of the model that works in guarded type theory.

Inspired by the solution of Giovannini et al., we refine the notion of abstract
categorical model of gradual typing into two novel classes of model suitable for
the guarded setting. We call these \emph{extensional} and \emph{intensional}
models.

% TODO more explanation

\subsection{Outline}

This work is organized as follows.

\begin{itemize}
  \item We begin in Section \ref{sec:idealized-models} with an overview of the
  existing abstract double-categorical semantics for gradual typing by New and
  Licata.
  \item We review the difficulties in adapting the New-Licata denotational model
  to the guarded setting.
  \item We motivate the refinement of the notion of abstract model into
  extensional and intensional models and present their definitions in full.
  \item We discuss the construction of an extensional model of gradual typing
  from a suitable ``base'' double category of effectful functions and relations.
  \item With the constructions in hand, we then instantiate the abstract
  framework to define concrete models of gradual typing capable of modeling
  languages with simple effects.
\end{itemize}

\subsection{Ambient Guarded Metalanguage}

Because guarded type theory plays a central role in the construction of the
concrete model in prior work, it is natural to expect that our abstract
framework should to some degree incorporate aspects of guarded type theory.
%
To that end, we demand the following of the categories involved in the
description of our abstract model:

\begin{enumerate}
    \item There is an applicative functor $\later \colon \calC \to \calC$.
    \item There is a morphism $\nxt \colon X \to \later X$ for any object $X$.
    \item For any morphism $f$ of the form $\later X \to X$, there is a unique
    morphism $\fix f \colon 1 \to X$ such that $\fix f = f \circ \nxt \circ \fix
    f$.
\end{enumerate}


% \section{Extensional and Intensional Models}

\section{Double Categorical Semantics of Graduality}\label{sec:idealized-models}

We begin with a review of abstract categorical models for gradual typing
described in prior work.
%
New and Licata \cite{new-licata18} modeled the graduality and type casts for
call-by-name gradual typing using \emph{double categories}, which are defined to
be categories internal to the category of categories. That is, a double category
$\mathcal C$ consists of a category $\mathcal C_f$ of ``objects and function
morphisms'' and a category $\mathcal C_{sq}$ of ``relation morphisms and
squares'' with functors (reflexive relation) $r : C_f \to C_{sq}$ and (source
and target) $s,t : C_{sq} \to C_f$ satisfying $sr = tr = \id$ as well as a
composition operation $c : C_{sq} \times_{s,t} C_{sq} \to C_{sq}$ respecting
source and target. This models an abstract notion of functions and relations.
For notation, we write function morphisms as $f : A \to B$ and relation
morphisms as $c : A \rel B$ where $c \in C_{sq}$ and $s(c) = A$ and $t(c) = B$.
Finally, a morphism $\alpha$ from $c$ to $d$ with $s(\alpha) = f$ and $s(\beta)
= g$ is visualized as

% https://q.uiver.app/#q=WzAsNCxbMCwwLCJBIl0sWzEsMCwiQiJdLFswLDEsIkEnIl0sWzEsMSwiQiciXSxbMCwyLCJmIiwyXSxbMSwzLCJnIl0sWzAsMSwiYyIsMCx7InN0eWxlIjp7ImJvZHkiOnsibmFtZSI6ImJhcnJlZCJ9LCJoZWFkIjp7Im5hbWUiOiJub25lIn19fV0sWzIsMywiZCIsMix7InN0eWxlIjp7ImJvZHkiOnsibmFtZSI6ImJhcnJlZCJ9LCJoZWFkIjp7Im5hbWUiOiJub25lIn19fV1d
\[\begin{tikzcd}[ampersand replacement=\&]
	A \& B \\
	{A'} \& {B'}
	\arrow["f"', from=1-1, to=2-1]
	\arrow["g", from=1-2, to=2-2]
	\arrow["c", "\shortmid"{marking}, no head, from=1-1, to=1-2]
	\arrow["d"', "\shortmid"{marking}, no head, from=2-1, to=2-2]
\end{tikzcd}\] 

Such a morphism is thought of as an abstraction of the notion of relatedness of
functions: functions take related inputs to related outputs. The composition
operations and functoriality give us a notion of composition of relations as
well as functions and vertical and horizontal composition of squares. In this
work we will be chiefly interested in \emph{locally thin} double categories,
that is, double categories where there is at most one square for any $f,c,g,d$.
In this case we use the notation $f \ltsq{c}{d} g$ to mean that a square like the
above exists.

New, Licata and Ahmed \cite{new-licata-ahmed2019} extended the axiomatic syntax
to call-by-push-value but did not analyze the structure categorically. We fill
in this missing analysis now, extending the double categorical semantics given
in \cite{new-licata18} from cartesian closed categories to CBPV models. 

We begin by reviewing the definition of a CBPV model:

\begin{definition}
  A CBPV model $\calM$ consists of:
  \begin{itemize}
    \item A cartesian-closed category $\calV$ of values and pure morphisms.
    \item A category $\calE$ of computations and linear morphisms.
    \item Functors $F : \calV \to \calE$ and $U : \calE \to \calV$ with $F \dashv U$.
    \item An action of $\calV$ on $\calE$, i.e., a functor $\arr \colon
    \calV^{op} \times \calE \to \calE$ satisfying the usual coherence laws,
    and such that $U(A \arr B)$ is isomorphic to $A \Rightarrow UB$.
  \end{itemize}

  A \emph{strict homomorphism of CBPV models} from $(\mathcal V,\mathcal
  E,\ldots)$ to $(\mathcal V', \mathcal E',\ldots)$ consists of functors $G_v :
  \mathcal V \to \mathcal V'$ and $G_e : \mathcal E \to \mathcal E'$ that
  strictly preserve all CBPV constructions; see the appendix for a more detailed
  definition. We call this a strict morphism in contrast to a \emph{lax}
  morphism, which only preserves CBPV constructions up to transformation.
\end{definition}

A model of the congruence rules of the CBPV system can be given by a locally
thin ``double CBPV model'', which we define below:

\begin{definition}\label{def:double-cbpv-model}
  A \emph{double CBPV model} is a category internal to the category of CBPV
  models and \emph{strict} homomorphisms of CBPV models\footnote{it may be
  possible to also define this as a notion of CBPV model internal to some
  structured $2$-category of categories, but we are not aware of any such
  definition of an internal CBPV model}.
\end{definition}

Part of the data of a double CBPV model can be visualized as follows:

% https://q.uiver.app/#q=WzAsNCxbMCwwLCJcXHZzcSJdLFsyLDAsIlxcZXNxIl0sWzAsMiwiXFx2ZiJdLFsyLDIsIlxcZWYiXSxbMiwzLCJcXEZmIiwwLHsiY3VydmUiOi0yfV0sWzMsMiwiXFxVZiIsMCx7ImN1cnZlIjotMn1dLFswLDEsIlxcRnNxIiwwLHsiY3VydmUiOi0yfV0sWzEsMCwiXFxVc3EiLDAseyJjdXJ2ZSI6LTJ9XSxbMiwwLCJcXHJ2Il0sWzAsMiwiXFx0diIsMCx7Im9mZnNldCI6LTEsImN1cnZlIjotMn1dLFswLDIsIlxcc3YiLDIseyJvZmZzZXQiOjEsImN1cnZlIjoyfV0sWzEsMywiXFx0ZSIsMCx7Im9mZnNldCI6LTEsImN1cnZlIjotMn1dLFsxLDMsIlxcc2UiLDIseyJvZmZzZXQiOjEsImN1cnZlIjoyfV0sWzMsMSwiXFxyZSJdLFs0LDUsIlxcYm90IiwxLHsic2hvcnRlbiI6eyJzb3VyY2UiOjIwLCJ0YXJnZXQiOjIwfSwic3R5bGUiOnsiYm9keSI6eyJuYW1lIjoibm9uZSJ9LCJoZWFkIjp7Im5hbWUiOiJub25lIn19fV0sWzYsNywiXFxib3QiLDEseyJzaG9ydGVuIjp7InNvdXJjZSI6MjAsInRhcmdldCI6MjB9LCJzdHlsZSI6eyJib2R5Ijp7Im5hbWUiOiJub25lIn0sImhlYWQiOnsibmFtZSI6Im5vbmUifX19XV0=
\[\begin{tikzcd}[ampersand replacement=\&]
	\vsq \&\& \esq \\
	\\
	\vf \&\& \ef
	\arrow[""{name=0, anchor=center, inner sep=0}, "\Fsq", curve={height=-12pt}, from=1-1, to=1-3]
	\arrow["\tv", shift left, curve={height=-12pt}, from=1-1, to=3-1]
	\arrow["\sv"', shift right, curve={height=12pt}, from=1-1, to=3-1]
	\arrow[""{name=1, anchor=center, inner sep=0}, "\Usq", curve={height=-12pt}, from=1-3, to=1-1]
	\arrow["\te", shift left, curve={height=-12pt}, from=1-3, to=3-3]
	\arrow["\se"', shift right, curve={height=12pt}, from=1-3, to=3-3]
	\arrow["\rv", from=3-1, to=1-1]
	\arrow[""{name=2, anchor=center, inner sep=0}, "\Ff", curve={height=-12pt}, from=3-1, to=3-3]
	\arrow["\re", from=3-3, to=1-3]
	\arrow[""{name=3, anchor=center, inner sep=0}, "\Uf", curve={height=-12pt}, from=3-3, to=3-1]
	\arrow["\bot"{description}, draw=none, from=0, to=1]
	\arrow["\bot"{description}, draw=none, from=2, to=3]
\end{tikzcd}\]

The diagram above makes it clear that there are two useful ways of viewing a
double CBPV model $\calM$. First, we can view $\calM$ as pair of CBPV models,
one forming the objects and vertical morphisms (the bottom of the diagram) and
the other forming the horizontal morphisms and the squares (the top of the
diagram). Alternatively, we can view $\calM$ as a pair of double categories, one
for the values (the left of the diagram), and one for the computations (the
right of the diagram). For the latter viewpoint, we will often write $\calC$ to
stand for an unspecified such double category associated with a double CBPV
model.

Given a double CBPV model, we can interpret the syntax of gradual typing as
follows. Type precision $A \ltdyn A'$ is interpreted as a relation morphism $c_A
: A \rel A'$ in $\mathcal V_{sq}$, and term precision $\Gamma \ltdyn \Gamma'
\vdash M \ltdyn M' : A \ltdyn A'$ is interpreted as a square $M
\ltsq{c_\Gamma}{UF c_A} M'$.
%
% TODO review this
The fact that $t,r$ and the composition are all given by strict CBPV
homomorphisms says that all the type constructors lift to precision
(monotonicity of type constructors) as well as all term constructors
(congruence). Further, $r$ and composition being strict homomorphisms implies
that all type constructors strictly preserve the identity relation (identity
extension) and composition.

Next, to model type casts, the New-Licata model further requires that every
value relation $c : A \rel A'$ is \emph{left representable} by a function $u_c :
A \to A'$ and every computation relation $d : B \rel B'$ is \emph{right
representable} by a function $d_c : B' \to B$. In a locally thin double
category, left and right representability are defined as follows:

\begin{definition} % TODO these are the simpler versions of the rules
  A relation $R : A \rel B$ is \emph{left representable} by $f : A \to B$ if $f
  \ltsq{R}{r(B)} \id$ and $\id \ltsq{r(A)}{R} f $.

  Dually, $R : A \rel B$ is \emph{right representable} by $g : B \to A$ if $\id
  \ltsq{R}{r(A)} g$ and $g \ltsq{r(B)}{R} \id$.
\end{definition}

% These squares say that the relation $c$ is \emph{representable} by the upcast
% morphism $\upc{c}$, essentially that the relation is a kind of \emph{graph} of
% the function in that $x \binrel{\sem c} y$ if and only if $\sem{\upc c}(x)
% \leq_{A_2} y$ \cite{shulman2007}.

These squares are sufficient to model the UpL/UpR/DnL/DnR rules for casts.
% TODO this is possible because of horizontal composition of squares
Additionally, since representable morphisms compose, the compositionality of
casts comes for free. However, the retraction property must be added as an
additional axiom to the model. 
%
To model the error being the least element, the New-Licata model adds the
requirement that $\mho \circ {!} \ltsq{r(A)}{r(UB)} f$ holds for all $f :
\mathcal V(A,B)$. Finally, the dynamic type can be modeled as an arbitrary value
type $D$ with arbitrary relations $\nat \rel D$ and $D \ra D \rel D$ and $D
\times D \rel D$ (or whatever basic type cases are required).


\section{Adapting the Notion of Model to Guarded Type Theory}


% we have seen in the previous section that there are obstructions to developing a
% concrete model in \emph{guarded} type theory. While guarded type theory makes
% the construction of arbitrary guarded recursive definitions possible, it comes
% at a significant cost to reasoning: unfolding of recursive definitions is
% explicit. In a denotational semantics, this means that non-well-founded
% recursion must perform a kind of \emph{observable step}. This presents a
% difficulty in proving graduality, which is a property that is oblivious to the
% number of steps that a program takes.

% Therefore, to prove graduality, we must be able to relate terms that take
% different numbers of steps. In the previous section we have examined two ways of
% doing this: in the first approach by defining an ordering $\semltbad$ that is
% closed under delays, and in the second by decomposing this ordering into a
% lock-step error ordering and a weak bisimilarity relation. But either way, we
% cannot escape the confines of Theorem \ref{thm:no-go}, which tells us that
% neither of these fomulations is transitive.

\subsection{Challenges in Adapting the New-Licata Model}

While the original double categorical semantics introduced in Section
\ref{sec:idealized-models} can be satisfied with classical domain theoretic
models as demonstrated by New and Licata \cite{new-licata18}, Giovannini et al.
\cite{gtt-sgdt} show that the New-Licata model cannot be directly adapted to
work in the guarded setting.

Key to understanding the issue is the fact that the equation for the dynamic
type involves a later, which means that when inspecting an element of the
dynamic type which is a function, an \emph{observable step} is required in the
semantics. This means that downcasts from the dynamic type may take an
observable step, and thus \emph{any} cast may take steps as it may involve a
downcast from the dynamic type functorially.

The issue in adapting the New-Licata model to the guarded setting is related to
the representability condition on the relations. At a high level, the problem
arises as a result of a tension between \emph{transitive reasoning} and
\emph{step-insensitivity}.
%
Transitive reasoning is desirable because it allows us to derive the cast rules
needed for graduality from the simpler, compositional versions of the rules.
This is crucial because the original rules quantify over an additional,
arbitrary type precision derivation and are therefore not syntax-independent.
%
On the other hand, the graduality property, as axiomatized by the
representability of type precision relations, is fundamentally a
step-insensitive property. The casts, which may involve observable steps, must
be related to the identity function, which does not step. The relation
modeling term precision must be able to account for this mismatch.

However, Giovannini et al. show that transitivity and unrestricted
step-insensitivity (plus a third minor condition) are incompatible as properties
of a relation in guarded type theory (see their Theorem 4.1 for a precise
statement).
%
The solution adopted by Giovannini et al. is to work with a \emph{lock-step
error ordering} that is transitive, and then restore a small amount of
step-insensitivity where needed through explicit synchronizing terms they call
\emph{perturbations}. The perturbations mimic the stepping behavior of casts,
but are otherwise equivalent to the identity, a property they formalize via a
notion of \emph{weak bisimilarity}. The authors relax the definition of
representability to allow for casts to be lock-step related to perturbations
rather than only to the identity function.
%
Finally, to define the denotation of term precision, the authors take the
closure of the lock-step ordering under weak bisimilarity on both sides:
%
\[ M \ltbisim N := \exists f, g.\, M \bisim f \ltdyn g \bisim N, \] 
%
where $\bisim$ is weak bisimilarity and $\ltdyn$ is the lock-step ordering.
%
The idea is that two terms are related if they can be synchronized in some
number of steps to two terms that are \emph{lock-step} related.
% This definition allows them to relate terms that differ up to a perturbation. 
% 
By separating term precision into a combination of weak bisimilarity and
lock-step reasoning, and using syntactic perturbations as a means to explicitly
synchronize terms as necessary, the authors achieve compositional reasoning
while still accommodating the steps taken by casts.


\subsection{Refining the Notion of Abstract Model}

As discussed in the previous section, the original notion of double categorical
model of graudal typing is not compatible with guarded type theory, where the
denotational semantics involves observable steps.
%
We are therefore led to \emph{refine} this notion of model to one that is
compatible with guarded type theory.
%
To motivate this refinement, consider the model of Giovannini et al.
\cite{gtt-sgdt} as reviewed in the previous section. Consider the following two
key observations about their model:
%
\begin{enumerate}
    \item The denotation of term precision (i.e., $\ltbisim$) is \emph{not}
    transitive. This follows from their no-go theorem (Theorem 4.1), as the
    relation satisfies the other two properties in the theorem and so it cannot
    be transitive.
    \item It is possible to derive the original representability rules for casts
    from the simplified versions using transitive reasoning.
\end{enumerate}
%
At first glance, these two observations seem contradictory. How can one use
transitivity if it does not hold? The key insight is that observations (1) and
(2) take place in two distinct settings. The derivation described in observation
(2) occurs \emph{before} the final denotation of term precision is defined. That
is, in (2), the relation in question is the lock-step ordering, \emph{not} its
closure under weak bisimilarity. The lock-step ordering is transitive, and thus
we can derive the required representability squares from the simplified ones via
horizontal composition of squares. Recall that perturbations are involved at
this point in order to keep terms in lock-step. Thus, much of the work in the
model construction consists of showing that the theory of perturbations is
compositional in the appropriate manner.

The existence of these two distinct settings in the concrete model motivates our
refinement of the idealized abstract model of gradual typing. In particular, we
divide the idealized notion of model into two classes, which we call
\emph{extensional} and \emph{intensional} models. The formal definitions will
follow, but the main idea can be understood in light of observations (1) and (2)
above. An \emph{extensional model} is an abstraction of the idea of working ``up
to weak bisimilarity'', i.e., extensionally. Such a model does not include
transitivity of term precision, since transitive reasoning and extensional
reasoning are incompatible in the guarded setting. On the other hand, an
extensional model does validate the rules for casts (in their original form).
% Such a model does not include transitivity of term precision, but does model
% the representability rules for casts (in their original form). 
Thus, we can interpret the syntax of gradual typing into an extensional model,
as the graduality property boils down to the validity of the representability
rules.

In contrast, an \emph{intensional model} is an abstraction of the idea of
working with a transitive, lock-step ordering, and using perturbations to
synchronize terms as needed. The restricted form of step-insensitivity provided
by perturbations allows us to reason up to observable steps without using weak
bisimilarity. Intensional models are appealing because they allow us to apply
compositional reasoning principles, e.g., deriving the original forms of the
cast rules from the simplified versions. On the other hand, intensional models
do not directly validate graduality, as we need to ``collapse'' the
perturbations to define the denotation of term precision. However, from an
intensional model, we can \emph{construct} an extensional model by performing
the abstract analogue of closing the lock-step ordering under weak bisimilarity.

In the remainder of this section, we define extensional and intensional models.
In the following section, we discuss how to construct an extensional model.
% ...and outline the steps to construct an extensional model.


% To this end, we develop an abstract notion of \emph{extensional} model where
% intuitively, ordering is up to weak bisimilarity. Transitivity of term precision
% does not hold in this setting, so the definition of extensional model weakens
% the idealized notion of model from double categories to reflexive graph
% categories. In particular, the definition drops the operation of horizontal
% pasting of squares, but it maintains a form of representability of relations
% crucial for graduality. An extensional model is sufficient for interpreting the
% terms and inequational theory of gradual typing. Unfortunately, as was mentioned
% in the previous section, constructing such a model directly is difficult: since
% we lack transitvity, we must take as primitive the original, more complicated
% representability rules. This makes the model dependent on the particular details
% of the cast calculus being modeled.

% However, the story does not end here. The approach of decomposing the ordering
% into the lock-step ordering and weak bisimilarity has the advantage that we can
% employ transitive reasoning when working with the lock-step ordering. The caveat
% is that we can no longer model casts as representable morphisms when reasoning
% with observable computation steps.
% %
% To axiomatize this situation abstractly, we introduce an intermediate notion of
% \emph{intensional} model, where intuitively, the ordering considers
% computational steps to be observable. Intensional models are based on double
% categories, but representability must be weakened to ``quasi-representability'',
% a kind of representability \emph{up to} observable steps. To reason up to
% observable steps without using weak bisimilarity, we develop a more refined
% notion we call \emph{perturbations}, certain terms that are bisimilar to the
% identity but can be manipulated explicitly in constructions. We then show that
% the compositional construction of casts from domain theoretic models can be
% adapted to this intensional setting by incorporating some explicit manipulation
% of perturbations. Finally, we show that taking an ``extensional collapse'' by
% weak bisimilarity provides a model that validates the graduality property.

\subsection{Extensional Models of Gradual Typing}

Since we lack transitivity of ordering when reasoning extensionally in guarded
type theory, our weakened notion of extensional model is based on reflexive
graph categories rather than double categories. This means we lose the reasoning
principle of horizontal pasting of squares. We will still require a notion of
\emph{composition} of relations, to model the transitivity of type precision. We
note that without horizontal pasting of squares, the notion of left/right
representability of relations is not sufficient to interpret the cast rules of
gradual typing. Instead, in an extensional model we work with a generalized
notion of representability that corresponds to the non-compositional form of the
cast rules:

\begin{definition}
    Let $c : A \rel A'$ and $f : A \to A'$ in a reflexive graph category with
    composition of relations. We say that $c$ is \emph{universally
    left-representable by} $f$ if for any $c_l : A_l \rel A$ and $c_r : A' \rel
    A_r$ we have $f \ltsq{cc_r}{c_r} \id$ and $\id \ltsq{c_l}{c_lc} f$.
    %
    Dually, let $d : B \rel B'$ and $\phi : B' \to B$.  We say that $d$ is
    \emph{universally right-representable by} $\phi$ if for any $d_l : B_l \rel
    B$ and $d_r : B' \rel B_r$ we have $\id \ltsq{d_ld}{d_l} \phi$ and $\phi
    \ltsq{d_r}{dd_r}\id$.
\end{definition}

In a reflexive graph category, these notions are stronger than the definitions
of left/right representability, since we can pick $c_l,c_r,d_l,d_r$ to be the
reflexive relations. In a double category these are equivalent, but the
equivalence uses horizontal pasting.

Additionally, while in the presence of horizontal pasting compositionality of
representable morphisms is automatic, without this principle we must require it
explicitly.
%
Finally, we explicitly add in a requirement that the type
constructors are functorial on relations up to order equivalence.

In summary:
\begin{definition}\label{def:extensional-model}
    An \emph{extensional model of gradual typing} consists of:

    \begin{enumerate}
    \item A locally thin reflexive graph internal to CBPV models.
    \item Composition of value and computation relations that form a category
    with the reflexive relations as identity. Call these categories $\mathcal
    V_r,\mathcal E_r$.
    \item Identity-on-objects functors $\upf : \mathcal V_r \to \mathcal V_f$
    and $\dnf : \mathcal E_r^{op} \to \mathcal E_f$ such that every $\upf\, c$
    universally left-represents $c$ and every $\dnf\, d$ universally
    right-represents $d$.
    \item The CBPV connectives $U,F,\times,\to$ are all \emph{covariant}
    functorial on relations up to equivalence: $U(dd') \equidyn U(d)U(d')$
    etc.\footnote{The reflexive graph structure already requires that these
    functors preserve identity relations.}
        \begin{itemize}
        \item $U(dd') \equidyn U(d)U(d')$
        \item $F(cc') \equidyn F(c)F(c')$
        \item $(cc') \to (dd') \equidyn (c \to d)(c' \to d')$
        \item $(c_1c_1') \times (c_2c_2') \equidyn (c_1 \times c_2)(c_1'\times c_2')$
        \end{itemize}
        where $c \equidyn c'$ means $\id \ltsq{c}{c'}\id$ and $\id \ltsq{c'}{c}
        \id$.
    \item A natural transformation $\mho : 1 \Rightarrow U$ such that $\mho
        \circ {!} \ltsq{r(A)}{r(UB)} f$ for any $f : A \to UB$.
    \item Distinguished value type $\nat$ with morphisms $z : \mathcal
        V(1,\nat)$ and $s : \mathcal V(\nat,\nat)$.
    \item Distinguished value types $D$ with distinguished relations $\iarr{}:
        U(D \to F D) \rel D$ and $\inat : \nat \rel D$ and $\itimes : D \times D
        \rel D$ each satisfying the retraction property $\dnc
        {\iarr{}}F(\upc{\iarr{}}) \equidyn \id$.
\end{enumerate}
\end{definition}

An extensional model of gradual typing is sufficient to interpret the syntax of
a gradual cast calculus.


\subsection{Intensional Models of Gradual Typing}

An intensional model of gradual typing is an abstraction of working in a setting
where the ordering on terms is transitive but requires terms to be in lock-step.
To aid in understanding the definition of an intensional model, we break the
definition into five steps, with each one building on the previous. The last
step will be the definition we have in mind when we refer to an intensional
model.

% The steps in the definition of intensional model are as follows:
% \begin{enumerate}
%   \item A double CBPV model
% \end{enumerate}

The starting point for an intensional model is the same as that of an idealized
model of gradual typing: namely, a double CBPV model (see Definition
\ref{def:double-cbpv-model}). In particular, unlike the extensional model where
we could only compose \emph{relations} horizontally, here we also have a
horizontal composition operation on squares.

% The starting point is similar to that of the extensional model. This time,
% however, since we are working intensionally, the semantic denotation of term
% precision \emph{is} transitive, so we \emph{do} have a horizontal composition
% operation on squares. Compare this to the extensional case, where we could only
% compose \emph{relations} horizontally, not squares.
% %
% We can specify this data compactly as a category internal to the category of
% CBPV models and lax morphisms, where we require that the reflexivity, source,
% and target morphisms are strict. In particular, as in the extensional case,
% there is a CBPV model of ``objects'' $\mathcal M_f$ and a CBPV model of
% ``arrows'' $\mathcal M_{sq}$. There are strict CBPV morphisms $r : \mathcal M_f
% \to \mathcal M_{sq}$ and $s, t : \mathcal M_{sq} \to \mathcal M_f$, just as
% before.
% %
% But now, we also have a CBPV morphism $m$ from the pullback $\mathcal M_{sq}
% \times_{s = t} \mathcal M_{sq}$ to $\mathcal M_{sq}$, i.e., ``composition of
% arrows''. In particular, this consists of a functor $m_{\mathcal{V}} : \vsq
% \times_{\sv = \tv} \vsq \to \vsq$ for composition of value relations/squares,
% and a functor $m_{\mathcal{E}} : \esq \times_{\se = \te} \esq \to \esq$ for
% composition of computation relations/squares. Furthermore, $s \circ m = s \circ
% \pi_1$ and $t \circ m = t \circ \pi_2$.

As in the definition of extensional model, we also require the existence of a
natural transformation $\mho : 1 \Rightarrow U$ such that $\mho\, \circ\,{!}
\ltsq{r(A)}{r(UB)} f$ for any $f : A \to UB$, and a distinguished value type
$\nat$ with morphisms $z : \mathcal V(1,\nat)$ and $s : \mathcal V(\nat,\nat)$.

To model intensional behavior, we specify that the model has a natural
transformation $\delta \colon U \Rightarrow U$. Specifically, this means there
is a family of vertical morphisms $\delta_B \in \vf(UB, UB)$, which is natural
in $B$ with respect to composition of vertical morphisms: for all $\phi \in
\ef(B_i, B_o)$ we have $U\phi \circ \delta_{B_1} = \delta_{B_2} \circ (U\phi)$.
Naturality of $\delta$ furthermore implies that for any $d : B \rel B'$, there
is a square $\delta_{B} \ltsq{Ud}{Ud} \delta_{B'}$.

We call a model satisfying these requirements an \emph{intensional pre-model}.
For the sake of ease of reference, we recap the definition below:
%
\begin{definition}
  An \textbf{intensional pre-model} consists of:
  \begin{itemize}
    \item A double CBPV model, i.e., a category internal to the category of CBPV
    models and lax morphisms, where we require that the morphisms $r$, $s$, and
    $t$ are strict.
    \item A natural transformation $\mho : 1 \Rightarrow U$ such that $\mho
    \circ {!} \ltsq{r(A)}{r(UB)} f$ for any $f : A \to UB$. 
    \item A value type $\nat$ with morphisms $z : \mathcal V(1,\nat)$ and $s :
    \mathcal V(\nat,\nat)$.
    \item A natural transformation $\delta \colon U \Rightarrow U$.
  \end{itemize}
\end{definition}


\begin{lemma}
  We have the following properties involving the delay maps in any intensional
  pre-model:
  \begin{enumerate}
    \item We can define a computation morphism $\delta_{FA}^* \in \ef(FA, FA)$ for
    all $A$. To do so, first compose the unit $\eta_A \in \vf(A, UFA)$ with the
    morphism $\delta_{FA} \in \vf(UFA, UFA)$, and then by the adjunction we get a
    computation morphism $\delta_{FA}^* \in \ef(FA, FA)$.
    \item There is a square $\delta_{FA}^* \ltdyn_{Fc}^{Fc} \delta_{FA'}^*$ for all $c : A \rel A$.
  \end{enumerate}
\end{lemma}
\begin{proof}
  See appendix.
\end{proof}

%   U\phi \circ \delta_{B_1}
% = U\phi \circ \theta_{B_1} \circ \nxt 
% = \theta_{B_2} \circ \later(U\phi) \circ \nxt
% = \theta_{B_2} \circ \nxt \circ U\phi
% = \delta_{B_2} \circ U\phi

\subsubsection{Weak Bisimilarity}\label{sec:abstract-model-bisimilarity}

The next ingredient in the definition of intensional model is \emph{weak
bisimilarity}, which formalizes the notion that two morphisms are the same ``up
to steps'', i.e., they differ only in that one may wait more than the other.
%
In particular, for any pair of objects $A$ and $A'$, in $\vf$, we require that
there is a reflexive, symmetric relation $\bisim_{A,A'}$ on the hom-set $\vf(A,
A')$, called the \emph{weak bisimilarity} relation. Similarly for the
computation category: there is a reflexive, symmetric relation $\bisim_{B,B'}$
defined on each hom-set $\ef(B, B')$.
%
Additionally, the weak bisimilarity relation should be a congruence with respect
to composition: if $f \bisim_{A,A'} f'$ and $g \bisim_{A',A''} g'$, then $g
\circ f \bisim_{A,A''} g' \circ f'$, and likewise for computations.

% It seems this is needed when defining the actions of the functors on semantic perturbations.
The bisimilarity relation should also be preserved by the functors $U, F, \arr$,
and $\times$ e.g., if $f \bisim_{A',A} g$ and $\phi \bisim_{B,B'} \phi'$, then
$f \arr \phi \bisim_{A \arr B, A' \arr B'} g \arr \phi'$. Similarly, the two
directions of the isomorphism of hom-sets $\vf(A, UB) \cong \ef(FA, B)$ should
preserve bisimilarity.

We can state all of this compactly by requiring that the categories $\vf$ and
$\ef$ be \emph{enriched} over the symmetric monoidal category $\setapprox$,
whose objects are sets equipped with a reflexive, symmetric relation and whose
arrows are functions between the sets that preserve the relation. Then we
require that the functors $U, F, \arr$, and $\times$ are
$\setapprox$-enriched, and likewise $F \dashv U$ as enriched functors.
%
Note that $\vsq$ and $\esq$ are not $\setapprox$-enriched, so in the functors
$r, s, t$ between $\vf$ and $\vsq$ and between $\ef$ and $\esq$, we view $\vf$
and $\ef$ as ordinary categories, where we simply forget the bisimilarity
relation on hom-sets.

Lastly, we require that for any computation object $B$, the delay map $\delta_B
\colon UB \to UB$ (which is defined in terms of $\theta_B$) is bisimilar to the
identity: $\delta_B \bisim_{UB, UB} \id_{UB}$.


\begin{definition}\label{def:step-1-model}   
An \emph{intensional pre-model with bisimilarity} is an intensional pre-model
$\calM$ such that:
\begin{itemize}
  % \item There is a reflexive, symmetric relation $\bisim$ on the hom-sets
  % $\vf(A,A')$ and $\ef(B,B')$ such that $\bisim$ preserves composition.
  \item The categories $\vf$ and $\ef$ are enriched over $\setapprox$.
  \item The functors $U, F, \arr$, and $\times$ are $\setapprox$-enriched, and
  the adjunction $F \dashv U$ is an enriched adjunction.
  \item For each $B$, the delay map $\delta_{B}$ is bisimilar to the identity
  $\delta_B \bisim_{UB, UB} \id_{UB}$.
\end{itemize}
\end{definition}

If $\calM$ is an intensional pre-model with bisimilarity, it will sometimes be
useful to have a concise definition describing the value objects, morphisms,
relations, and squares as a coherent unit, and likewise for the computations. To
do this, we introduce some terminology:

\begin{definition}\label{def:bisim-dbl-cat}
  A \emph{bisimilarity reflexive-graph} $\calC$ consists of:
  \begin{itemize}
    \item A category $\cf$ of objects and vertical morphisms, enriched in
    $\setapprox$.
    \item A thin category $\csq$ of horizontal morphisms and squares.
    \item A functor $r: |\cf| \to \csq$ where $|\cf|$ is the underlying ordinary
    category of $\cf$.
    \item Functors $s, t: \csq \to |\cf|$ such that $sr = tr = \id$.
  \end{itemize}

  A \emph{functor} $G : \calC \to \calC'$ of bisimilarity reflexive-graphs
  consists of:
  \begin{itemize}
    \item A $\setapprox$-enriched functor $G_f: \cf \to \cf'$.
    \item A functor $G_{sq} : \csq \to \csq'$.
    \item Source and target are preserved: $s' \circ G_{sq} = G_f \circ s$, and $t' \circ G_{sq} = G_f \circ t$.
    \item Reflexivity is preserved: $r' \circ G_f = G_{sq} \circ r$.
  \end{itemize}

  Let $G, H: \calC \to \calC'$ be functors of bisimilarity reflexive-graphs. A
  \emph{natural transformation} $\alpha: G \to H$ consists of:
  \begin{itemize}
    \item A $\setapprox$-enriched natural transformation $\alpha^f: G_f \to H_f$.
    \item A natural transformation $\alpha^{sq}: G_{sq} \to H_{sq}$.
  \end{itemize}

  A \emph{bisimilarity double-category} is a bisimilarity reflexive graph that
  additionally has a ``horizontal composition'' functor $\circ \colon \csq
  \times_{s = t} \csq \to \csq$ where $\csq \times_{s = t} \csq$ is the
  following pullback:

    % https://q.uiver.app/#q=WzAsNCxbMCwwLCJcXGNzcSBcXHRpbWVzX3tzID0gdH0gXFxjc3EiXSxbMCwxLCJcXGNzcSJdLFsxLDAsIlxcY3NxIl0sWzEsMSwiXFxjZiJdLFswLDFdLFsxLDNdLFsyLDNdLFswLDJdLFswLDMsIiIsMSx7Im9mZnNldCI6Miwic3R5bGUiOnsibmFtZSI6ImNvcm5lciJ9fV1d
  \[\begin{tikzcd}[ampersand replacement=\&]
    {\csq \times_{s = t} \csq} \& \csq \\
    \csq \& \cf
    \arrow[from=1-1, to=1-2]
    \arrow[from=1-1, to=2-1]
    \arrow["\lrcorner"{anchor=center, pos=0.125}, shift right=2, draw=none, from=1-1, to=2-2]
    \arrow[from=1-2, to=2-2]
    \arrow[from=2-1, to=2-2]
  \end{tikzcd}\]

  We additionally require that composition is well-behaved with respect to the
  source and target functors, and satisfies associativity and left and right
  unit laws, as in an ordinary double category.

% % https://q.uiver.app/#q=WzAsNSxbMSwwLCJcXGNzcSBcXHRpbWVzX3tzID0gdH0gXFxjc3EiXSxbMSwyLCJcXGNmIl0sWzEsMSwiXFxjc3EiXSxbMCwxLCJcXGNzcSJdLFsyLDEsIlxcY3NxIl0sWzAsMiwiXFxjaXJjIl0sWzIsMSwicyIsMix7Im9mZnNldCI6MX1dLFszLDEsInMiLDJdLFswLDMsIlxccGlfMSIsMl0sWzAsNCwiXFxwaV8yIl0sWzQsMSwidCJdLFsyLDEsInQiLDAseyJvZmZzZXQiOi0xfV1d
% \[\begin{tikzcd}[ampersand replacement=\&]
% 	\& {\csq \times_{s = t} \csq} \\
% 	\csq \& \csq \& \csq \\
% 	\& \cf
% 	\arrow["{\pi_1}"', from=1-2, to=2-1]
% 	\arrow["\circ", from=1-2, to=2-2]
% 	\arrow["{\pi_2}", from=1-2, to=2-3]
% 	\arrow["s"', from=2-1, to=3-2]
% 	\arrow["s"', shift right, from=2-2, to=3-2]
% 	\arrow["t", shift left, from=2-2, to=3-2]
% 	\arrow["t", from=2-3, to=3-2]
% \end{tikzcd}\]

  A \emph{(lax) functor of bisimilarity double-categories} is defined
  analogously to a functor of bisimilarity reflexive-graphs, except we
  additionally require that the functor preserves composition of horizontal
  morphisms up to a square:

%   % https://q.uiver.app/#q=WzAsNSxbMCwwLCJHWF8xIl0sWzEsMCwiR1hfMiJdLFsyLDAsIkdYXzMiXSxbMCwxLCJHWF8xIl0sWzIsMSwiR1hfMyJdLFszLDQsIkcoUiBcXGNpcmMgUicpIiwyLHsic3R5bGUiOnsiYm9keSI6eyJuYW1lIjoiYmFycmVkIn0sImhlYWQiOnsibmFtZSI6Im5vbmUifX19XSxbMCwxLCJHUiIsMCx7InN0eWxlIjp7ImJvZHkiOnsibmFtZSI6ImJhcnJlZCJ9LCJoZWFkIjp7Im5hbWUiOiJub25lIn19fV0sWzEsMiwiR1InIiwwLHsic3R5bGUiOnsiYm9keSI6eyJuYW1lIjoiYmFycmVkIn0sImhlYWQiOnsibmFtZSI6Im5vbmUifX19XSxbMCwzLCJcXGlkIiwyXSxbMiw0LCJcXGlkIl1d
% \[\begin{tikzcd}[ampersand replacement=\&]
% 	{GX_1} \& {GX_2} \& {GX_3} \\
% 	{GX_1} \&\& {GX_3}
% 	\arrow["GR", "\shortmid"{marking}, no head, from=1-1, to=1-2]
% 	\arrow["\id"', from=1-1, to=2-1]
% 	\arrow["{GR'}", "\shortmid"{marking}, no head, from=1-2, to=1-3]
% 	\arrow["\id", from=1-3, to=2-3]
% 	\arrow["{G(R \circ R')}"', "\shortmid"{marking}, no head, from=2-1, to=2-3]
% \end{tikzcd}\]

% https://q.uiver.app/#q=WzAsNSxbMCwwLCJHWF8xIl0sWzEsMCwiR1hfMiJdLFsyLDAsIkdYXzMiXSxbMCwxLCJHWF8xIl0sWzIsMSwiR1hfMyJdLFszLDQsIkcoUiBcXGNpcmMgUicpIiwyLHsic3R5bGUiOnsiYm9keSI6eyJuYW1lIjoiYmFycmVkIn0sImhlYWQiOnsibmFtZSI6Im5vbmUifX19XSxbMCwxLCJHUiIsMCx7InN0eWxlIjp7ImJvZHkiOnsibmFtZSI6ImJhcnJlZCJ9LCJoZWFkIjp7Im5hbWUiOiJub25lIn19fV0sWzEsMiwiR1InIiwwLHsic3R5bGUiOnsiYm9keSI6eyJuYW1lIjoiYmFycmVkIn0sImhlYWQiOnsibmFtZSI6Im5vbmUifX19XSxbMiw0LCIiLDAseyJvZmZzZXQiOi0xLCJzdHlsZSI6eyJoZWFkIjp7Im5hbWUiOiJub25lIn19fV0sWzAsMywiIiwyLHsib2Zmc2V0IjoxLCJzdHlsZSI6eyJoZWFkIjp7Im5hbWUiOiJub25lIn19fV0sWzAsMywiIiwxLHsic3R5bGUiOnsiaGVhZCI6eyJuYW1lIjoibm9uZSJ9fX1dLFsyLDQsIiIsMSx7InN0eWxlIjp7ImhlYWQiOnsibmFtZSI6Im5vbmUifX19XV0=
\[\begin{tikzcd}[ampersand replacement=\&]
	{GX_1} \& {GX_2} \& {GX_3} \\
	{GX_1} \&\& {GX_3}
	\arrow["GR", "\shortmid"{marking}, no head, from=1-1, to=1-2]
	\arrow[shift right, no head, from=1-1, to=2-1]
	\arrow[no head, from=1-1, to=2-1]
	\arrow["{GR'}", "\shortmid"{marking}, no head, from=1-2, to=1-3]
	\arrow[shift left, no head, from=1-3, to=2-3]
	\arrow[no head, from=1-3, to=2-3]
	\arrow["{G(R \circ R')}"', "\shortmid"{marking}, no head, from=2-1, to=2-3]
\end{tikzcd}\]

\end{definition}

Bisimilarity double-categories, functors between them, and natural
transformations of such functors form a 2-category, which we denote
$\bisimdblcat$:
%
\begin{lemma}
  There is a 2-category $\bisimdblcat$ whose objects are bisimilarity
  double-categories, whose 1-cells are functors between bisimilarity
  double-categories, and whose 2-cells are natural transformations between such
  functors.
\end{lemma}
\begin{proof}
  See appendix.
\end{proof}

% TODO:
% Opposite of a bisim double-cat 
% Product of bisim double-cats
% Exponential of bisim double-cats

With these definitions in hand, we can give an alternative definition of an
intensional pre-model with bisimilarity:

\begin{definition}
  An intensional pre-model with bisimilarity consists of the following diagram
  in $\bisimdblcat$:

  % https://q.uiver.app/#q=WzAsMixbMCwwLCJcXGNhbFYiXSxbMiwwLCJcXGNhbEUiXSxbMCwxLCJGIiwwLHsiY3VydmUiOi00fV0sWzEsMCwiVSIsMV0sWzEsMCwiVSIsMSx7ImN1cnZlIjotNH1dLFsxLDAsIlxcRGVsdGEoMV9cXGNhbFYpIiwwLHsib2Zmc2V0IjotMywiY3VydmUiOi01fV0sWzIsMywiIiwwLHsibGV2ZWwiOjEsInN0eWxlIjp7Im5hbWUiOiJhZGp1bmN0aW9uIn19XSxbMyw0LCJcXGRlbHRhIiwwLHsic2hvcnRlbiI6eyJzb3VyY2UiOjIwLCJ0YXJnZXQiOjIwfX1dLFs1LDQsIlxcbWhvIiwyLHsic2hvcnRlbiI6eyJzb3VyY2UiOjIwLCJ0YXJnZXQiOjIwfX1dXQ==
\[\begin{tikzcd}[ampersand replacement=\&]
	\calV \&\& \calE
	\arrow[""{name=0, anchor=center, inner sep=0}, "F", curve={height=-24pt}, from=1-1, to=1-3]
	\arrow[""{name=1, anchor=center, inner sep=0}, "U"{description}, from=1-3, to=1-1]
	\arrow[""{name=2, anchor=center, inner sep=0}, "U"{description}, curve={height=-24pt}, from=1-3, to=1-1]
	\arrow[""{name=3, anchor=center, inner sep=0}, "{\Delta(1_\calV)}", shift left=3, curve={height=-30pt}, from=1-3, to=1-1]
	\arrow["\dashv"{anchor=center, rotate=-90}, draw=none, from=0, to=1]
	\arrow["\delta", shorten <=3pt, shorten >=3pt, Rightarrow, from=1, to=2]
	\arrow["\mho"', shorten <=2pt, shorten >=2pt, Rightarrow, from=3, to=2]
\end{tikzcd}\]

  Along with a functor of bisimilarity double-categories $\arr \colon \calV^{op}
  \times \calE \to \calE$ satisfying $U(A \arr B) \cong A \Rightarrow UB$
  
  Which additionally satisfies:
  \begin{itemize}
    \item 
  \end{itemize}
\end{definition}


\subsubsection{Perturbations}\label{sec:abstract-model-perturbations}

We next focus on the specification of perturbations in a pre-model with
bisimilarity. To begin, let $\calC$ be a bisimilarity double-category.
%
% Let $\calM$ be a pre-model with bisimilarity, and let $\calC$ stand
% for the double category of values of $\calM$, or the double category of
% computations of $\calM$.
%
For an object $X$ of $\calC$, a \emph{semantic perturbation} of $X$ is an
endomorphism $f \in \cf(X, X)$ such that $f \bisim_{X,X} \id_X$. The collection
of semantic perturbations on $X$ forms a monoid under composition, with the
identity element being the identity $\id_X$. We write $\semptb{X}$ to refer to
this monoid.

For every object $X \in \cf$, we require that there is a monoid $M_X$ of
\emph{syntactic perturbations} of $X$ and a monoid homomorphism $i_X : M_X \to
\semptb{X}$ that interprets these as semantic perturbations.

The prototypical semantic perturbation for a computation object $B$ is the delay
map $\delta_B \colon UB \to UB$. We require that for any computation object $B$,
the monoid $M_{UB}$ contains an element $m$ such that $i_{UB}(m) = \delta_B$.

We require that the functors $\times, \timesk$, $\arr, \tok$, $U, \Uk$, $F, \Fk$
all preserve perturbations. As in the prior work of Giovannini et al., we also
require that the perturbations interact nicely with the relation morphisms.
Namely, if $R : X \rel X'$ is a relation morphism, we require there to be two
homomorphisms of monoids $\push_R : M_X \to M_{X'}$ and $\pull_R : M_{X'} \to
M_X$. The authors call this the \emph{push-pull property}.

If $\calM$ is a pre-model with bisimilarity that additionally satisfies the
above properties, we say that $M$ is a \emph{pre-model with perturbations}. To
summarize:
%
% TODO
\begin{definition}\label{def:step-2-model} 
  A \emph{pre-model with perturbations} $\calM$ is a pre-model with
  bisimilarity, such that the following additional requirements are met:
  \begin{enumerate}
    \item For each object (value or computation) $X$, there is a monoid $M_X$
    and a homomorphism of monoids $i_X : M_X \to \semptb{X}$.
    \item For all $B$, there is an element $m$ such that $i_{UB}(m) = \delta_B$.
    \item The functors $\times, \timesk$, $\arr, \tok$, $U, \Uk$, $F, \Fk$
    preserve perturbations.
    \item The push-pull property holds for all $c : A \rel A'$ and all $d : B
    \rel B'$.
  \end{enumerate}
\end{definition}


\subsection{Representable Relations}\label{sec:abstract-model-representable}

Now we turn to the notion of representability of horizontal morphisms, which
allows us to specify the behavior of casts.
%
As in the case of an idealized abstract model, we want that each value
horizontal morphism $c$ is left-representable by an embedding morphism $e_c$ and
$Fc$ is right-representable by a projection morphism $p_c$, and dually for
computation horizontal morphisms $d$. However, since casts take observable
steps, we must insert perturbations on the side opposite the cast to keep both
sides in lock-step. This is because the squares in an intensional model are
intended to model a strong error ordering between vertical morphisms in which
both sides must have the same stepping behavior. Following Giovannini et al., we
call this weakened notion of representability up-to an explicit perturbation
\emph{quasi-representability}.

In this section, we begin with the notion of intensional pre-model with
perturbations as defined in the previous section, and provide the additional
conditions that axiomatize the behavior of casts. The precise definitions are as
follows.

First, let $\mathcal M$ be any double category with a notion of perturbations,
i.e., for any object $X$ in $\mathcal M$ there is a monoid $\Ptb_X$ with a monoid
homomorphism into the endomorphisms on $X$.
%
\begin{definition}\label{def:quasi-left-representable}
  Let $R$ be a horizontal morphism in $\mathcal M$ between objects $X$ and $Y$.
  We say that $R$ is \emph{quasi-left-representable by} a vertical morphism $f$
  in $\mathcal M(X, Y)$ if there are perturbations $\delle_R \in \Ptb_X$ and
  $\delre_R \in \Ptb_Y$ such that there is a square $\upl : f \ltsq{R}{r(Y)}
  \ptb(\delre_R)$ and a square $\upr : \ptb(\delle_R) \ltsq{r(X)}{R} f$.

  \begin{center}
    \begin{tabular}{ m{7em} m{7em} } 
      % UpL
      \begin{tikzcd}[ampersand replacement=\&]
        X \& {Y} \\
        {Y} \& {Y}
        \arrow["f"', from=1-1, to=2-1]
        \arrow["\ptb(\delta_R^{r,e})", from=1-2, to=2-2]
        \arrow["R", "\shortmid"{marking}, no head, from=1-1, to=1-2]
        \arrow["r(Y)", "\shortmid"{marking}, no head, from=2-1, to=2-2]
      \end{tikzcd}
      &
      % UpR
      \begin{tikzcd}[ampersand replacement=\&]
        X \& {X} \\
        {X} \& {Y}
        \arrow["\ptb(\delta_R^{l,e})"', from=1-1, to=2-1]
        \arrow["f", from=1-2, to=2-2]
        \arrow["r(X)", "\shortmid"{marking}, no head, from=1-1, to=1-2]
        \arrow["R"', "\shortmid"{marking}, no head, from=2-1, to=2-2]
      \end{tikzcd}
    \end{tabular}
  \end{center}
%
\end{definition}

\begin{definition}\label{def:quasi-right-representable}
  Let $R$ be a horizontal morphism between $X$ and $Y$. We say that $R$ is
  \emph{quasi-right-representable by} $f \in \mathcal M(Y, X)$ if there exist
  perturbations $\dellp_R \in \Ptb_X$ and $\delrp_R \in \Ptb_Y$ such that we
  have a square $\dnr : \ptb(\dellp_R) \ltsq{R}{r(X)} f$ and a square $\dnl : f
  \ltsq{r(Y)}{R} \ptb(\delrp_R)$.
  %
  % \begin{center}
  %   \begin{tabular}{ m{7em} m{7em} } 
  %     % DnR
  %     \begin{tikzcd}[ampersand replacement=\&]
  %       {X} \& {Y} \\
  %       {X} \& {X}
  %       \arrow["\delta_R^{l,p}"', from=1-1, to=2-1]
  %       \arrow["f", from=1-2, to=2-2]
  %       \arrow["R", "\shortmid"{marking}, no head, from=1-1, to=1-2]
  %       \arrow[from=2-1, to=2-2, Rightarrow, no head]
  %     \end{tikzcd}
  %     &
  %     % DnL
  %     \begin{tikzcd}[ampersand replacement=\&]
  %       {Y} \& {Y} \\
  %       {X} \& {Y}
  %       \arrow["f"', from=1-1, to=2-1]
  %       \arrow["\delta_R^{r,p}", from=1-2, to=2-2]
  %       \arrow[from=1-1, to=1-2, Rightarrow, no head]
  %       \arrow["R"', "\shortmid"{marking}, no head, from=2-1, to=2-2]
  %     \end{tikzcd}
  %   \end{tabular}
  % \end{center}
  % %
  % We call the first square $\dnr$ and the second square $\dnl$.
  %
\end{definition}

We can now define a notion of an intensional model where the horizontal
morphisms are quasi-representable:
%
\begin{definition}\label{def:step-3-model}
  Let $\calM$ be an intensional pre-model with perturbations. We say that
  $\calM$ is a \emph{pre-model with representability} if:
  \begin{itemize}
    \item Every value horizontal morphism $c : A \rel A'$ is
    quasi-left-representable in the double category $\calV$. That is, there is a
    morphism $\upc{c} \in \calV(A, A')$ that quasi-left-represents $c$.
    \item Every computation horizontal morphism $d : B \rel B'$ is
    quasi-right-representable in $\calE$. That is, there is a morphism $\dnc{d}
    \in \calE(B', B)$ that quasi-right-represents $d$.
  \end{itemize}
\end{definition}

\subsection{The Dynamic Type}\label{sec:abstract-model-dyn}

We now discuss how to specify the denotation of the dynamic type in an
intensional model. It is important to note that it makes sense to axiomatize the
dynamic type already in a base intensional pre-mdoel. That is, we do not require
the notions of weak bisimilarity, perturbations, or representability of
horizontal morphisms to define the denotation of the dynamic type.
% However, for our purposes, we will assume we are given an intensional
% pre-model with representability.

% Given such a model $\calM$, to interpret the dynamic type we require that
% there is...

Let $\calM$ be an intensional pre-model, a pre-model with bisimilarity, a
pre-model with perturbations, or a pre-model with representability. To interpret
the dynamic type, we require that there is a distinguished value object $D \in
\calV$, such that for every value object $A$, there is a distinguished
horizontal morphism $\inj{A} \colon A \rel D$.

Now we can state the definition of an intensional model:
\begin{definition}
  An \emph{intensional model of gradual typing} is an intensional pre-model with
  representability that furthermore has an interpretation of the dynamic type.
\end{definition}

\section{Constructing a Guarded Model of Gradual Typing}

In this section, we discuss how to \emph{construct} an abstract model of gradual
typing via a step-by-step process. Recall that the ultimate goal is to build an
extensional model of gradual typing (Definition \ref{def:extensional-model}), as
such a model contains the structure necessary to interpret the syntax and
precision rules of a gradual cast calculus.

Although the definition of an extensional model involves a significant amount of
data, it is not necessary to create an instance of an extensional model from
scratch. Much of the difficult work can be carried out abstractly. In
particular, it is possible to construct an extensional model from an intensional
model using the abstract analogue of the bisimilarity-closure construction in
the work of Giovannini et al. Moreover, much of the construction of an
intensional model can also be carried out abstractly.

% \subsection{Starting Point: Intensional Pre-Model with Bisimilarity}

We take as a starting point an intensional pre-model with bisimilarity
(Definition \ref{def:step-1-model}).

\subsection{Adding Perturbations}

Suppose we are given an intensional pre-model with bisimilarity $\calM$. Our
goal is to construct a model $\calM'$ with perturbations, as in Definition
\ref{def:step-2-model}.

The summary of the construction is as follows:
%
\begin{itemize}
  \item
  A value object in $\calM'$ consists of a value object $A$ of $\calM$
  paired with a monoid $\Ptb_A$ of syntactic perturbations, and a homomorphism of
  monoids $\ptb_A : \Ptb_A \to \semptb{A}$ interpreting the perturbations as
  semantic perturbations.
  %
  Dually, a computation object of $\calM'$ consists of a computation object $B$ of
  $\calM$ paired with a monoid $\Ptb_B$, along with a monoid homomorphism $\ptb_B
  : \Ptb_B \to \semptb{B}$.

  \item
  The value and computation vertical morphisms in $\calM'$ are defined simply as
  those of $\calM$; they do not interact with the perturbations.

  \item 
  % TODO is this a property of the relation, or additional structure?
  A value relation $c: (A, \Ptb_A, \ptb_A) \rel (A', \Ptb_{A'}, \ptb_{A'})$ in
  $\calM'$ is defined to be a value relation $c : A \rel A'$ in $\calM$ that
  satisfies the push-pull property with respect to the monoids $\Ptb_A$ and
  $\Ptb_{A'}$.
  %
  Dually, a computation relation $d: (B, \Ptb_B, \ptb_B) \rel (B', \Ptb_{B'},
  \ptb_{B'})$ in $\calM'$ is defined to be a computation relation in $\calM$ that
  satisfies the push-pull property with respect to $\Ptb_B$ and $\Ptb_{B'}$.

  \item
  Lastly, value and computation squares in $\calM'$ are defined to be the same as
  those in $\calM$.
\end{itemize}

\subsection{Adding Representability}

\subsection{The Dynamic Type}

Lastly, suppose we are given an intensional pre-model with representability
$\calM$, along with an object $D$ with distinguished relations $\inat, \itimes,
\iarr$.
%
We construct a model for which each value object $A$ has a corresponding
horizontal morphism $\inj{c} : A \rel D$.

\subsection{Extracting an Extensional Model: The Bisimilarity-Closure Construction}

Lastly, we can obtain an extensional model from an intensional model using the
abstract analogue of the bisimilarity-closure definition in the work of
Giovannini et al. In particular, let $\calM$ be an intensional model. We define
an extensional model $\calM'$ as follows. The objects and vertical/horizontal
morphisms are the same as those in $\calM$; it is only the definition of squares
that differs.

Specifically, let $f : X_i \to X_o$ and $g : X_i' \to X_o'$ be vertical
morphisms (of either value or computation types). Let $R_i : X_i \rel X_i'$ and
$R_o : X_o \rel X_o'$ be horizontal morphisms. We define
%
\[ f \ltsqbisim{R_i}{R_o} g := 
  \exists f', g'.\, f \bisim_{X_i, X_o} f' \ltsq{R_i}{R_o} g' \bisim_{X_i', X_o'} g. \]
%
It remains to verify that the model satisfies the requirements of an extensional
model. 
%
% TODO deriving the original cast rules from the simplified ones in the intensional model
%
The basic idea is that the perturbations in the intensional squares are all
bisimilar to the identity function, and hence they disappear when we perform the
above construction.

\section{Constructions on Abstract Models}

In this section, we explain how to augment a model of gradual typing with
additional effects presented in the form of an adjunction on the value or
computation categories.
%
We assume the minimum amount of data given, i.e., an intensional pre-model with
bisimilarity $\calM$ (see Definition \ref{def:step-1-model}), along with a
suitable value object interpreting the dynamic type.

%
% We assume here that the categories modeling the value and computation types are
% set-based (as opposed to a Kripke model of presheaves). The main obligation will
% be to construct the double categories of values and computations and the pair of
% adjoint functors $F$ and $U$ between these double categories.
%
% Importantly, we also assume that the equation for the dynamic type is of the
% same form as described in the previous section. Since the construction of the
% perturbations for the dynamic type does not depend on the specific
% implementations of the functors $U$, $F$, $\times$ and $\arr$, we need only
% construct the underlying type, the ordering, and the bisimilarity relation for
% the dynamic type.

% Our approach is as follows. Suppose we are given an intensional pre-model with
% bisimilarity $\calM$ (see Definition \ref{def:step-1-model}).

Recall that we can visualize (part of) the underlying double CBPV model of
$\calM$ as follows:
%
% https://q.uiver.app/#q=WzAsNCxbMCwwLCJcXHZzcSJdLFsyLDAsIlxcZXNxIl0sWzAsMiwiXFx2ZiJdLFsyLDIsIlxcZWYiXSxbMiwzLCJcXEZmIiwwLHsiY3VydmUiOi0yfV0sWzMsMiwiXFxVZiIsMCx7ImN1cnZlIjotMn1dLFswLDEsIlxcRnNxIiwwLHsiY3VydmUiOi0yfV0sWzEsMCwiXFxVc3EiLDAseyJjdXJ2ZSI6LTJ9XSxbMiwwLCJcXHJ2Il0sWzAsMiwiXFx0diIsMCx7Im9mZnNldCI6LTEsImN1cnZlIjotMn1dLFswLDIsIlxcc3YiLDIseyJvZmZzZXQiOjEsImN1cnZlIjoyfV0sWzEsMywiXFx0ZSIsMCx7Im9mZnNldCI6LTEsImN1cnZlIjotMn1dLFsxLDMsIlxcc2UiLDIseyJvZmZzZXQiOjEsImN1cnZlIjoyfV0sWzMsMSwiXFxyZSJdLFs0LDUsIlxcYm90IiwxLHsic2hvcnRlbiI6eyJzb3VyY2UiOjIwLCJ0YXJnZXQiOjIwfSwic3R5bGUiOnsiYm9keSI6eyJuYW1lIjoibm9uZSJ9LCJoZWFkIjp7Im5hbWUiOiJub25lIn19fV0sWzYsNywiXFxib3QiLDEseyJzaG9ydGVuIjp7InNvdXJjZSI6MjAsInRhcmdldCI6MjB9LCJzdHlsZSI6eyJib2R5Ijp7Im5hbWUiOiJub25lIn0sImhlYWQiOnsibmFtZSI6Im5vbmUifX19XV0=
\[\begin{tikzcd}[ampersand replacement=\&]
	\vsq \&\& \esq \\
	\\
	\vf \&\& \ef
	\arrow[""{name=0, anchor=center, inner sep=0}, "\Fsq", curve={height=-12pt}, from=1-1, to=1-3]
	\arrow["\tv", shift left, curve={height=-12pt}, from=1-1, to=3-1]
	\arrow["\sv"', shift right, curve={height=12pt}, from=1-1, to=3-1]
	\arrow[""{name=1, anchor=center, inner sep=0}, "\Usq", curve={height=-12pt}, from=1-3, to=1-1]
	\arrow["\te", shift left, curve={height=-12pt}, from=1-3, to=3-3]
	\arrow["\se"', shift right, curve={height=12pt}, from=1-3, to=3-3]
	\arrow["\rv", from=3-1, to=1-1]
	\arrow[""{name=2, anchor=center, inner sep=0}, "\Ff", curve={height=-12pt}, from=3-1, to=3-3]
	\arrow["\re", from=3-3, to=1-3]
	\arrow[""{name=3, anchor=center, inner sep=0}, "\Uf", curve={height=-12pt}, from=3-3, to=3-1]
	\arrow["\bot"{description}, draw=none, from=0, to=1]
	\arrow["\bot"{description}, draw=none, from=2, to=3]
\end{tikzcd}\]
%
Now suppose we are given a pair of adjunctions $\Ff' \dashv \Uf'$ and $\Fsq'
\dashv \Usq'$, of the form
%
% https://q.uiver.app/#q=WzAsNCxbMCwwLCJcXHZzcSJdLFsyLDAsIlxcdnNxIl0sWzAsMiwiXFx2ZiJdLFsyLDIsIlxcdmYiXSxbMiwzLCJcXEZmJyIsMCx7ImN1cnZlIjotMn1dLFszLDIsIlxcVWYnIiwwLHsiY3VydmUiOi0yfV0sWzAsMSwiXFxGc3EnIiwwLHsiY3VydmUiOi0yfV0sWzEsMCwiXFxVc3EnIiwwLHsiY3VydmUiOi0yfV0sWzIsMCwiXFxydiJdLFswLDIsIlxcc3YiLDAseyJvZmZzZXQiOi0xLCJjdXJ2ZSI6LTJ9XSxbMCwyLCJcXHR2IiwyLHsib2Zmc2V0IjoxLCJjdXJ2ZSI6Mn1dLFsxLDMsIlxcdHYiLDAseyJvZmZzZXQiOi0xLCJjdXJ2ZSI6LTJ9XSxbMSwzLCJcXHN2IiwyLHsib2Zmc2V0IjoxLCJjdXJ2ZSI6Mn1dLFszLDEsIlxccnYiXSxbNCw1LCJcXGJvdCIsMSx7InNob3J0ZW4iOnsic291cmNlIjoyMCwidGFyZ2V0IjoyMH0sInN0eWxlIjp7ImJvZHkiOnsibmFtZSI6Im5vbmUifSwiaGVhZCI6eyJuYW1lIjoibm9uZSJ9fX1dLFs2LDcsIlxcYm90IiwxLHsic2hvcnRlbiI6eyJzb3VyY2UiOjIwLCJ0YXJnZXQiOjIwfSwic3R5bGUiOnsiYm9keSI6eyJuYW1lIjoibm9uZSJ9LCJoZWFkIjp7Im5hbWUiOiJub25lIn19fV1d
\[\begin{tikzcd}[ampersand replacement=\&]
	\vsq \&\& \vsq \\
	\\
	\vf \&\& \vf
	\arrow[""{name=0, anchor=center, inner sep=0}, "{\Fsq'}", curve={height=-12pt}, from=1-1, to=1-3]
	\arrow["\sv", shift left, curve={height=-12pt}, from=1-1, to=3-1]
	\arrow["\tv"', shift right, curve={height=12pt}, from=1-1, to=3-1]
	\arrow[""{name=1, anchor=center, inner sep=0}, "{\Usq'}", curve={height=-12pt}, from=1-3, to=1-1]
	\arrow["\tv", shift left, curve={height=-12pt}, from=1-3, to=3-3]
	\arrow["\sv"', shift right, curve={height=12pt}, from=1-3, to=3-3]
	\arrow["\rv", from=3-1, to=1-1]
	\arrow[""{name=2, anchor=center, inner sep=0}, "{\Ff'}", curve={height=-12pt}, from=3-1, to=3-3]
	\arrow["\rv", from=3-3, to=1-3]
	\arrow[""{name=3, anchor=center, inner sep=0}, "{\Uf'}", curve={height=-12pt}, from=3-3, to=3-1]
	\arrow["\bot"{description}, draw=none, from=0, to=1]
	\arrow["\bot"{description}, draw=none, from=2, to=3]
\end{tikzcd}\]
%
We seek conditions under which the compositions of these adjunctions forms a
double CBPV model with bisimilarity of the form
%
% https://q.uiver.app/#q=WzAsNixbMCwwLCJcXHZzcSJdLFsyLDAsIlxcdnNxIl0sWzAsMiwiXFx2ZiJdLFsyLDIsIlxcdmYiXSxbNCwwLCJcXGVzcSJdLFs0LDIsIlxcZWYiXSxbMiwzLCJcXEZmJyIsMCx7ImN1cnZlIjotMn1dLFszLDIsIlxcVWYnIiwwLHsiY3VydmUiOi0yfV0sWzAsMSwiXFxGc3EnIiwwLHsiY3VydmUiOi0yfV0sWzEsMCwiXFxVc3EnIiwwLHsiY3VydmUiOi0yfV0sWzIsMCwiXFxydiJdLFswLDIsIlxcc3YiLDAseyJvZmZzZXQiOi0xLCJjdXJ2ZSI6LTJ9XSxbMCwyLCJcXHR2IiwyLHsib2Zmc2V0IjoxLCJjdXJ2ZSI6Mn1dLFsxLDMsIlxcdHYiLDAseyJvZmZzZXQiOi0xLCJjdXJ2ZSI6LTJ9XSxbMSwzLCJcXHN2IiwyLHsib2Zmc2V0IjoxLCJjdXJ2ZSI6Mn1dLFszLDEsIlxccnYiXSxbMSw0LCJcXEZzcSIsMCx7ImN1cnZlIjotMn1dLFs0LDEsIlxcVXNxIiwwLHsiY3VydmUiOi0yfV0sWzMsNSwiXFxGZiIsMCx7ImN1cnZlIjotMn1dLFs1LDMsIlxcVWYiLDAseyJjdXJ2ZSI6LTJ9XSxbNCw1LCJcXHRlIiwyLHsib2Zmc2V0IjotMSwiY3VydmUiOi0yfV0sWzQsNSwiXFxzZSIsMix7Im9mZnNldCI6MSwiY3VydmUiOjJ9XSxbNSw0LCJcXHJlIl0sWzYsNywiXFxib3QiLDEseyJzaG9ydGVuIjp7InNvdXJjZSI6MjAsInRhcmdldCI6MjB9LCJzdHlsZSI6eyJib2R5Ijp7Im5hbWUiOiJub25lIn0sImhlYWQiOnsibmFtZSI6Im5vbmUifX19XSxbOCw5LCJcXGJvdCIsMSx7InNob3J0ZW4iOnsic291cmNlIjoyMCwidGFyZ2V0IjoyMH0sInN0eWxlIjp7ImJvZHkiOnsibmFtZSI6Im5vbmUifSwiaGVhZCI6eyJuYW1lIjoibm9uZSJ9fX1dLFsxNiwxNywiIiwwLHsibGV2ZWwiOjEsInN0eWxlIjp7Im5hbWUiOiJhZGp1bmN0aW9uIn19XSxbMTgsMTksIiIsMCx7ImxldmVsIjoxLCJzdHlsZSI6eyJuYW1lIjoiYWRqdW5jdGlvbiJ9fV1d
\[\begin{tikzcd}[ampersand replacement=\&]
	\vsq \&\& \vsq \&\& \esq \\
	\\
	\vf \&\& \vf \&\& \ef
	\arrow[""{name=0, anchor=center, inner sep=0}, "{\Fsq'}", curve={height=-12pt}, from=1-1, to=1-3]
	\arrow["\sv", shift left, curve={height=-12pt}, from=1-1, to=3-1]
	\arrow["\tv"', shift right, curve={height=12pt}, from=1-1, to=3-1]
	\arrow[""{name=1, anchor=center, inner sep=0}, "{\Usq'}", curve={height=-12pt}, from=1-3, to=1-1]
	\arrow[""{name=2, anchor=center, inner sep=0}, "\Fsq", curve={height=-12pt}, from=1-3, to=1-5]
	\arrow["\tv", shift left, curve={height=-12pt}, from=1-3, to=3-3]
	\arrow["\sv"', shift right, curve={height=12pt}, from=1-3, to=3-3]
	\arrow[""{name=3, anchor=center, inner sep=0}, "\Usq", curve={height=-12pt}, from=1-5, to=1-3]
	\arrow["\te"', shift left, curve={height=-12pt}, from=1-5, to=3-5]
	\arrow["\se"', shift right, curve={height=12pt}, from=1-5, to=3-5]
	\arrow["\rv", from=3-1, to=1-1]
	\arrow[""{name=4, anchor=center, inner sep=0}, "{\Ff'}", curve={height=-12pt}, from=3-1, to=3-3]
	\arrow["\rv", from=3-3, to=1-3]
	\arrow[""{name=5, anchor=center, inner sep=0}, "{\Uf'}", curve={height=-12pt}, from=3-3, to=3-1]
	\arrow[""{name=6, anchor=center, inner sep=0}, "\Ff", curve={height=-12pt}, from=3-3, to=3-5]
	\arrow["\re", from=3-5, to=1-5]
	\arrow[""{name=7, anchor=center, inner sep=0}, "\Uf", curve={height=-12pt}, from=3-5, to=3-3]
	\arrow["\bot"{description}, draw=none, from=0, to=1]
	\arrow["\dashv"{anchor=center, rotate=-90}, draw=none, from=2, to=3]
	\arrow["\bot"{description}, draw=none, from=4, to=5]
	\arrow["\dashv"{anchor=center, rotate=-90}, draw=none, from=6, to=7]
\end{tikzcd}\]





\section{Instantiating the Abstract Framework}

Now that we have demonstrated how to construct a guarded model of gradual
typing, we instantiate our framework with several examples of concrete
gradually-typed languages.

\subsection{A Model for the Simple Gradually-Typed Lambda Calculus}\label{sec:gtlc-concrete-model}

We begin with the gradually-typed lambda calculus as described in Section 2 of
the paper of Giovannini et al. \cite{gtt-sgdt}. Following their construction, we
define an intensional pre-model with bisimilarity for this language as follows:

\begin{itemize}
  \item Value objects are \emph{predomains} --- sets $A$ equipped with a partial
  ordering relation $\ltdyn_A$ and a reflexive, symmetric relation $\bisim_A$.

  \item Value vertical morphisms are \emph{predomain morphisms}. A predomain
  morphism between $A_i$ and $A_o$ is a function $f : A_i \to A_o$ between the
  underlying sets that preserves the ordering and bisimilarity relations.

  \item Value horizontal morphisms are \emph{predomain relations}. A predomain
  relation $c$ between $A$ and $A'$ is a relation between the underlying sets
  that is closed under the predomain ordering relations on both sides.
  Composition of predomain relations is given by the usual relational
  composition.

  \item A value square specifies the usual logical relations relationship
  between predomain morphisms: $f \ltsq{c_i}{c_o} g$ means that for any $x
  \binrel{c} y$, we have $f(x) \binrel{c_o} g(y)$.

  \item Computation objects are \emph{error domains} --- predomains $B$ that
  additionally have a least element $\mho_B \ltdyn b$, and a morphism of
  predomains $\theta: \later B \to B$ such that $\delta_B := \theta_B \circ
  \nxt$ is bisimilar to the identity morphism.

  \item Computation vertical morphisms are \emph{error domain morphisms}
  
  \item Computation horizontal morphisms are \emph{error domain relations}. An
  error domain relation $d : B \rel B'$ consists of a relation between the
  underlying predomains that additionally respects the error and $\theta$
  elements.
  %
  Composition of error domain relations $d$ and $d'$ is not given by composition
  of the underlying predomain relations; instead, it is defined inductively as
  the smallest error domain relation containing the composition of the
  underlying predomain relations.
  
  \item A computation square $\phi \ltsq{d}{d'} \phi'$ is defined analogously to
  a value square.
\end{itemize}

The functor $U$ from error domains to predomains simply forgets the error
element and the $\theta$ map.
%
The functor $F$ from predomains to error domains takes a predomain $A$ to the
\emph{free error domain} on $A$. The underlying predomain of $FA$ is defined as
the unique solution to the guarded recursive type equation
%
\[ U(FA) \cong A + 1 + \later U(FA), \]
%
where we name the constructors $\eta: A \to U(FA)$, $\mho: UFA$, and $\theta:
\later U(FA) \to U(FA)$.

The ordering relation on $FA$ is the \emph{lock-step} error ordering defined in
Figure \ref{fig:lock-step-error-ordering}, which also defines the weak
bisimilarity relation on $FA$.

\begin{figure}
  \begin{minipage}{0.5\textwidth}
    \fbox{$l_1 \ltls l_2$}
    \begin{align*}
        &\eta\, x \ltls \eta\, y \text{ if } 
            x \mathbin{\ltdyn_A} y \\
        %		
        &\mho \ltls l' \\
        %
        &\theta\, \tilde{l} \ltls \theta\, \tilde{l'} \text{ if } 
            \later_t (\tilde{l}_t \ltls \tilde{l'}_t)
    \end{align*}
  \end{minipage}
  \begin{minipage}{0.5\textwidth}
  \fbox{$l_1 \bisim l_2$}
    \begin{align*}
        \eta\, x \bisim \eta\, y &\text{ iff } x \bisim_A y \\
      %
        \mho \bisim \mho &\text{ iff } \top \\
      %		
        \theta\, \tilde{x} \bisim \theta\, \tilde{y} &\text{ iff } \later_t (\tilde{x}_t \bisim \tilde{y}_t) \\
      %	
        \theta\, \tilde{x} \bisim \mho &\text{ iff } \exists n. \theta\, \tilde{x} = \delta^n(\mho)\\
      %	
        \theta\, \tilde{x} \bisim \eta\, y &\text{ iff } \exists n. \exists x \bisim_A y.
          (\theta\, \tilde{x} = \delta^n(\eta\, x))\\
      %
        \mho \bisim \theta\, \tilde{y} &\text{ iff } \exists n. \theta\, \tilde{y} = \delta^n(\mho) \\
      %	
        \eta\, x \bisim \theta\, \tilde{y} &\text { iff } \exists n. \exists y \bisim_A x. (\theta\, \tilde{y} = \delta^n (\eta\, y))
      \end{align*}
  \end{minipage}
    \caption{The lock-step error ordering and weak bisimilarity relations on $FA$.}
    \label{fig:lock-step-error-ordering}
\end{figure}

Given two predomains $A$ and $A'$, we can form their product $A \times A'$,
where the partial ordering and bisimilarity relations are the obvious extensions
to the product of the underlying sets.
%
Given predomains $A_i$ and $A_o$, the set of morphisms of predomains from $A_i$
to $A_o$ also forms a predomain, denoted $A_i \Rightarrow A_o$, with the
ordering and bisimilarity relations defined in the obvious way, e.g., $f \bisim
g := \forall x \bisim y.\, f(x) \bisim g(y)$.

In particular, when the codomain is of the form $U(B)$ for some error domain
$B$, then the predomain $A \Rightarrow UB$ inherits the structure of an error
domain from that of $B$. The error element is the constant function returning
$\mho_B$, and the $\theta$ map takes a later-morphism $\tilde{f}$ and returns
the morphism $\lambda x. \theta_B((\later \tilde{f})\, x)$ using the action of
$\later$ on morphisms.

It is straightforward to verify that the functors $\times, \arr, U,$ and $F$
extend to relations and squares.


Lastly, to define the denotation of the dynamic type, we first construct a
predomain $|D|$ satisfying
%
\[ |D| \cong \mathbb{N} + |D| \times |D| + U(D \arr FD), \]
%
where the functors in question are understood to be operating on the categories
of predomains and error domains as appropriate. In particular, $\later$ acts on
predomains: for a predomain $A$, the ordering and bisimilarity relations on
$\later A$ are defined in the obvious manner --- two elements of $\later A$ are
related if one time step later, they are related in $A$.
%
Having defined the predomain for the dyanmic type, we then define the monoid of
perturbations for the dynamic type by least-fixpoint in the category of monoids
as
\[ 
  M_D \cong (M_D \oplus M_D) \oplus 
  (\mathbb{N} \oplus M_D^{op} \oplus \mathbb{N} \oplus M_D),
\]
and the interpretation $i_D$ of the perturbations as endomorphisms of $|D|$ is
defined straightfowardly. The resulting triple of $(|D|, M_D, i_D)$ is taken as
the denotation of the dynamic type.

% The composition $UF$ defines a (strong) monad on the category of predomains.

% \subsection{Additional Set-Based Examples}

% We next consider additional examples where the value types are set-based (as
% opposed to a Kripke model of presheaves). The main obligation will be to
% construct the double categories of values and computations and the pair of
% adjoint functors $F$ and $U$ between these double categories.
% %
% Importantly, we assume that the equation for the dynamic type is the same. Since
% the construction of the perturbations for the dynamic type does not depend on
% the specific implementations of the functors $U$, $F$, $\times$ and $\arr$, we
% need only construct the underlying type, ordering relation, and bisimilarity
% relation for the dynamic type.

\begin{comment}
\subsection{New Models From Old}

%\subsubsection{Global State}


We begin by considering \emph{global state} as an effect. Recall that, given an
adjunction $F \dashv U$ with $F : C \to D$, we can augment $C$ with global state
for a state object $S \in C$ by composing it with the usual product-exponential
adjunction: % TODO citation?
%
% https://q.uiver.app/#q=WzAsMyxbMCwwLCJDIl0sWzEsMCwiQyJdLFsyLDAsIkQiXSxbMCwxLCJTIFxcdGltZXMgLSIsMCx7ImN1cnZlIjotMX1dLFsxLDAsIlMgXFxSaWdodGFycm93IC0iLDAseyJjdXJ2ZSI6LTF9XSxbMSwyLCJGIiwwLHsiY3VydmUiOi0xfV0sWzIsMSwiVSIsMCx7ImN1cnZlIjotMX1dLFszLDQsIiIsMCx7ImxldmVsIjoxLCJzdHlsZSI6eyJuYW1lIjoiYWRqdW5jdGlvbiJ9fV0sWzUsNiwiIiwwLHsibGV2ZWwiOjEsInN0eWxlIjp7Im5hbWUiOiJhZGp1bmN0aW9uIn19XV0=
\[\begin{tikzcd}[ampersand replacement=\&]
	C \& C \& D
	\arrow[""{name=0, anchor=center, inner sep=0}, "{S \times -}", curve={height=-6pt}, from=1-1, to=1-2]
	\arrow[""{name=1, anchor=center, inner sep=0}, "{S \Rightarrow -}", curve={height=-6pt}, from=1-2, to=1-1]
	\arrow[""{name=2, anchor=center, inner sep=0}, "F", curve={height=-6pt}, from=1-2, to=1-3]
	\arrow[""{name=3, anchor=center, inner sep=0}, "U", curve={height=-6pt}, from=1-3, to=1-2]
	\arrow["\dashv"{anchor=center, rotate=-90}, draw=none, from=0, to=1]
	\arrow["\dashv"{anchor=center, rotate=-90}, draw=none, from=2, to=3]
\end{tikzcd}\]


% https://q.uiver.app/#q=WzAsMyxbMCwwLCJcXGNhbFZfZiJdLFsxLDAsIlxcY2FsVl9mIl0sWzIsMCwiXFxjYWxFX2YiXSxbMCwxLCJTIFxcdGltZXMgLSIsMCx7ImN1cnZlIjotMX1dLFsxLDAsIlMgXFxSaWdodGFycm93IC0iLDAseyJjdXJ2ZSI6LTF9XSxbMSwyLCJGIiwwLHsiY3VydmUiOi0xfV0sWzIsMSwiVSIsMCx7ImN1cnZlIjotMX1dLFszLDQsIiIsMCx7ImxldmVsIjoxLCJzdHlsZSI6eyJuYW1lIjoiYWRqdW5jdGlvbiJ9fV0sWzUsNiwiIiwwLHsibGV2ZWwiOjEsInN0eWxlIjp7Im5hbWUiOiJhZGp1bmN0aW9uIn19XV0=
\[\begin{tikzcd}[ampersand replacement=\&]
	{\calV_f} \& {\calV_f} \& {\calE_f}
	\arrow[""{name=0, anchor=center, inner sep=0}, "{S \times -}", curve={height=-6pt}, from=1-1, to=1-2]
	\arrow[""{name=1, anchor=center, inner sep=0}, "{S \Rightarrow -}", curve={height=-6pt}, from=1-2, to=1-1]
	\arrow[""{name=2, anchor=center, inner sep=0}, "F", curve={height=-6pt}, from=1-2, to=1-3]
	\arrow[""{name=3, anchor=center, inner sep=0}, "U", curve={height=-6pt}, from=1-3, to=1-2]
	\arrow["\dashv"{anchor=center, rotate=-90}, draw=none, from=0, to=1]
	\arrow["\dashv"{anchor=center, rotate=-90}, draw=none, from=2, to=3]
\end{tikzcd}\]

\end{comment}

% \subsubsection{Generalizing Error Domains: Guarded Interaction Trees}

\subsection{New Models From Scratch}

Next, we discuss how we can build a new model of gradual typing from scratch
given a collection of effects presented algebraically. Our goal will be to
generalize the notion of error domain to accommodate these effects.

% We now explore how to generalize the notion of error domain and the free error
% domain construction to accommodate additional effects besides errors and
% computational steps. We do this in two phases: first, we work at the level of
% sets and functions, and then we consider the ordering and bisimilarity relations.

% Recall that the input to the abstract model construction is an intensional
% pre-model with bisimilarity, in addition to a suitable object modeling the
% dynamic type.

More concretely, let $\Sigma: \PreDom \to \PreDom$ be an endofunctor on the
category of predomains (see the previous section for the definition of
predomain). The purpose of $\Sigma$ is to represent a signature of effects.
First, to model gradual typing we need to ensure that signature has cases for
errors and computational steps. To do this systematically, we define
$\Sigmahat := \Sigma \uplus \Delta(1) \uplus \text{Id}$, where $\uplus$ is
the coproduct of functors, and $\Delta(1)$ is the constant functor returning the
terminal predomain $1$.
%
We will then be interested in algebras of $\Sigmahat \circ \later$. Here, we
note that this $\later$ is operating on predomains. Furthermore, this definition
implies that all occurrences of the input $X$ are guarded by a $\later$. A
further generalization would be to consider algebras of a functor where some
occurrences of $X$ are unguarded and others are guarded. In this case, we would
need to assume more about our given functor $\Sigma$, e.g., that it is presented
as a \emph{polynomial} endofunctor; we will pursue this in the subsequent
section.

More concretely, an algebra of the functor $\Sigmahat \circ \later$ consists of
a predomain $|B|$ along with a morphism of predomains
%
$\alg_B \colon \Sigmahat (\later |B|) \to |B|$.
%
Note that $\Sigmahat (\later |B|) = \Sigma(\later |B|) \uplus 1 \uplus (\later
|B|)$. Thus, an algebra map $\alg_B$ induces three morphisms of predomains; we
denote these by $\oper_B \colon \Sigma(\later |B|) \to |B|$, $\mho_B \colon 1
\to |B|$, and $\theta_B \colon \later |B| \to |B|$.
%
Notice also that, by definition of morphism of predomains, these maps preserve
the ordering and bisimilarity relations of the relevant predomains.


% We write $\op_B$ to refer to the map $\alg_B \circ \inone \colon \Sigma(\later B) \to B$,
% where $\inone$ is the first injection into the coproduct (recall the definition of $\Sigmahat$).
% We write $\mho_B$ to refer to the map $\alg_B \circ \intwo \colon 1 \to B$,
% We write $\theta_B$ to refer to the map $\alg_B \circ \inthree \colon \later B \to B$.

We can now define the object that will serve as the analogue of error domains in
our construction:
%
\begin{definition}\label{def:sigma-domain}
  Let $\Sigma: \PreDom \to \PreDom$ be a functor. A \emph{$\Sigma$-domain} is an
  algebra of the functor $\Sigmahat \circ \later$. 
  %i.e., a predomain $B$ along with a morphism of predomains $\alg_B \colon
  %\Sigmahat (\later B) \to B$.
  Furthermore, the element $\mho_B \colon 1 \to B$ has the property that $\mho_B
  \ltdyn_B x$ for all $x \in B$. % TODO does this follow from an algebraic equation?
\end{definition}

\begin{example}
  For $\Sigma := \Delta(0)$, where $0$ is the empty predomain, we recover the
  original notion of error domain from Section \ref{sec:gtlc-concrete-model}.
\end{example}

% TODO define the action of later on predomain relations?
As with ordinary error domains, there are corresponding notions of morphisms,
relations, and squares of $\Sigma$-domains:
%
\begin{definition}\label{def:sigma-domain-mor-rel-sq}
  Fix $\Sigma$ as above, and let $B$ and $B'$ be $\Sigma$-domains. A
  \emph{morphism of $\Sigma$-domains} from $B$ to $B$' is a morphism $\phi$ of
  the underlying predomains that preserves the algebraic structure, i.e.,
  $\alg_{B'} \circ (\Sigmahat \circ \later)(\phi) = \phi \circ \alg_B$.
  %
  Spelling this out concretely, we have in particular that
  %
  (1) $\phi \circ \oper_B = \oper_{B'} \circ ((\Sigmahat \circ \later)(\phi))$,
  %
  (2) $\phi(\mho_B) = \mho_{B'}$, and
  %
  (3) $\phi \circ \theta_B = \theta_{B'} \circ (\later \phi)$.

  A \emph{relation of $\Sigma$-domains} $d : B \rel B'$ is a relation $|d|$ of
  the underlying predomains that is a congruence with respect to the algebraic
  structure, i.e., if $x \binrel{\Sigmahat(\later |d|)} y$, then $\alg_B(x)
  \binrel{|d|} \alg_B(y)$. (Here, we have used the actions of the predomain
  functor $\Sigmahat$ and the predomain functor $\later$ on predomain
  relations.) 
  %
  Spelling out the above requirement, we have in particular
  %
  (1) If $x \binrel{(\Sigma(\later |d|))} y$, then $\oper_B(x) \binrel{|d|} \oper_{B'}(y)$.
  %
  (2) $\mho_B \binrel{|d|} \mho_{B'}$, and
  %
  (3) If $\later_t(\tilde{x}_t \binrel{|d|} \tilde{y}_t)$, then
  $\theta_B(\tilde{x}) \binrel{|d|} \theta_{B'}(\tilde{y})$

  A \emph{square of $\Sigma$-domains} $\phi \ltsq{d}{d'} \phi'$ is defined
  analogously to a square of predomains. 
\end{definition}

\begin{remark}
  The requirement that $\mho_B$ be related to $\mho_{B'}$ in the definition of
  relation may seem weaker than the corresponding requirement in the definition
  of a relation of error domains in the previous section, where we required that
  $\mho_B \binrel{|d|} y$ for \emph{all} $y \in B'$. However, the former implies
  the latter using the fact that predomain relations are closed under the
  ordering on both sides and that $\mho_{B'} \ltdyn y$ for all $y$.
\end{remark}

% TODO composition of \Sigma-domain relations

Now we proceed to describe an intensional pre-model with bisimilarity based on
an adjunction between predomains and $\Sigma$-domains. More specifically:
%
\begin{itemize}
  \item As in Section \ref{sec:gtlc-concrete-model}, the value objects, vertical
  morphisms, horizontal morphisms, and squares are respectively predomains,
  morphisms of predomains, predomain relations, and predomain squares.

  \item The computation objects are $\Sigma$-domains as defined above. Vertical
  morphisms, horizontal morphisms, and squares are respectively morphisms,
  relations, and squares of $\Sigma$-domains, as defined above.
\end{itemize}

Observe that the predomain $A \To UB$ consisting of predomain morphisms $A \to
UB$ inherits the $\Sigma$-domain structure of $B$.

Next, we define the adjunction $F \dashv U$ between predomains and
$\Sigma$-domains. The functor $U$ simply forgets the additional algebraic
structure. The functor $F$ constructs, for a predomain $A$, the free
$\Sigma$-domain on $A$. The underlying set is defined by guarded recursion as
the type with the following constructors:
%
\begin{align*}
  \GT{\Sigma}{A} &:=      \GTret  \colon A \to \GT{\Sigma}{A} \\
              &\mid \GTop         \colon \Sigma(\laterhs \GT{\Sigma}{A}) \to \GT{\Sigma}{A}
              &\mid \mho_\Sigma   \colon \GT{\Sigma}{A} \\
              &\mid \theta_\Sigma \colon \laterhs \GT{\Sigma}{A} \to \GT{\Sigma}{A} \\
\end{align*}
%
The ordering and bisimilarity relations on $\GT{\Sigma}{A}$ are a generalization
of the lock-step error ordering and weak bisimilarity from the previous section.
The full definitions are given in Figure \ref{fig:gtree-ordering-bisim}. Notice
in the cases when both sides are a $\GTop$, we pass the relation being defined
to $\Sigma$, but under a later.
%
Formally, since $\Sigma$ is a functor on \emph{predomains}, what we are actually
doing is applying the guarded-recursive assumption that a predomain structure
for the set $\GT{\Sigma}{A}$ exists later. Alternatively, instead of applying
$\Sigma$ to a predomain, one can think of this as simply applying its action on
the ordering or bisimilarity relations to the ordering/bisimilarity relation
that we assume exists later.

\begin{figure}
  \begin{minipage}{0.5\textwidth}
    \fbox{$l_1 \ltls l_2$}
    \begin{align*}
        &\GTret\, x \ltls \GTret\, y \text{ if } x \mathbin{\ltdyn_A} y \\
        %
        &\GTop\, x \ltls \GTop\, y \text{ if } x \mathbin{\Sigma(\laterhs \ltls)} y \\
        %
        &\mho_\Sigma \ltls l' \\
        %
        &\theta\, \tilde{l} \ltls \theta\, \tilde{l'} \text{ if } 
            \later_t (\tilde{l}_t \ltls \tilde{l'}_t)
    \end{align*}
  \end{minipage}
  \begin{minipage}{0.5\textwidth}
    \fbox{$l_1 \bisim l_2$}
    \begin{align*}
        \GTret\, x \bisim \GTret\, y &\text{ iff } x \bisim_A y \\
      %		
        \mho \bisim \mho &\text{ iff } \top \\
      %
        \theta\, \tilde{x} \bisim \theta\, \tilde{y} &\text{ iff } \later_t (\tilde{x}_t \bisim \tilde{y}_t) \\
      %	
        \GTop\, x \bisim \GTop\, y \text{ if } x \mathbin{\Sigma(\laterhs \bisim)} y \\
      %
      %
      %
        \theta\, \tilde{x} \bisim \mho &\text{ iff } \exists n. \theta\, \tilde{x} = \delta^n(\mho)\\
      %	
        \theta\, \tilde{x} \bisim \eta\, y &\text{ iff } \exists n. \exists x \bisim_A y.
          (\theta\, \tilde{x} = \delta^n(\eta\, x))\\
      %
        \mho \bisim \theta\, \tilde{y} &\text{ iff } \exists n. \theta\, \tilde{y} = \delta^n(\mho) \\
      %	
        \eta\, x \bisim \theta\, \tilde{y} &\text { iff } \exists n. \exists y \bisim_A x. (\theta\, \tilde{y} = \delta^n (\eta\, y))
      \end{align*}
  \end{minipage}
    \caption{The lock-step error ordering and weak bisimilarity relations on $\GT{\Sigma}{A}$.}
    \label{fig:gtree-ordering-bisim}
\end{figure}

We observe that $\GT{\Sigma}{A}$ is the free $\Sigma$-domain on $A$: the set of
morphisms of predomains $A \to UB$, where $B$ is a $\Sigma$-domain, is in
bijection with morphisms of $\Sigma$-domains $\GT{\Sigma}{A} \to B$.

Additionally, we define the action of the functor $F$ on predomain relations $c
: A \rel A'$ as the heterogeneous version of the above lock-step ordering.

\subsection{Generalization: Allowing for Non-Guarded Occurrences}

In the previous section, we required that all recursive occurrences of the
$\Sigma$-domain were guarded by a later. We can generalize this situation to
allow for non-guarded occurrences, provided they are positive.

\subsection{Example: State}




\begin{comment}

Suppose we are given a container $E = (S \colon \text{Type}, P \colon S \to
\text{Type})$. Intuitively, $S$ represents the names of a collection of effects,
and for each $s : S$, the arity of the effect is given by $Ps$. For a container
$E$, the polynomial endofunctor presented by $E$ is defined by $P_E(X) =
\Sigma_{s : S} \Pi_{p : P s} X$.
%
To model gradual typing, we need to ensure that our containers have cases for
errors and computational steps. To do this systematically, given a container $E
= (S, P)$ as above, we will define a \emph{modified container} $\hat{E}$ by
$(S', P')$ where $S' := S + 2$, and $P'$ is defined as follows. We define $P'
(\inl\, s) = P s$, and $P' (\inr\, \text{true}) = \bot$, and $P' (\inr\,
\text{false}) = \top$.
%
We will then be interested in algebras of the functor $P_{\hat{E}} \circ
\later$. Note that this means all recursive occurrences are guarded by a
$\later$. A further generalization would be to consider the case where some
occurrences are inductive and others are guarded-recursive; we do not pursue
this approach here.

\begin{definition}
  Let $E$ be a container. An $E$-domain is an algebra of the functor
  $P_{\hat{E}} \circ \later$, i.e., a set $B$ along with a map $a \colon
  P_{\hat{E}} (\later B) \to B$.
\end{definition}

\begin{example}
  When $E$ is the empty container, we recover the usual notion of error domains.
\end{example}


To define the free-forgetful adjunction, we use the notion of \emph{guarded
interaction tree} \cite{guarded-interaction-trees}. We recall the basics of
guarded interaction trees below; see their paper for more details.
%
Given a ground type $A$ and a signature of effects $E$, we define the type of
guarded interaction trees for $A$ and $E$, denoted $\GT{E}{A}$, via guarded
recursion as the datatype with the following constructors:
%
\begin{align*}
  \GT{E}{A} &:=   \GTret   \colon A \to \GT{E}{A} \\
            &\mid \mho_E   \colon \GT{E}{A} \\
            &\mid \theta_E \colon \laterhs \GT{E}{A} \to \GT{E}{A} \\
            &\mid \GTeff   \colon \Pi_{i \in I}(
                                    \ins_i(\laterhs \GT{E}{A})
                                    \times (\outs_i(\laterhs \GT{E}{A}) \to \laterhs \GT{E}{A})
                                  ) \to \GT{E}{A}
\end{align*}
%
Here, the effect signature $E$ consists of a set of operation names $I$ and for
each $i \in I$, functions $\ins_i$ and $\outs_i$ from Type to Type that
determine the arities of the operations. The $\GTret$ constructor takes an
element $x$ of $A$ and returns a tree that 

\end{comment}



\section{Discussion}

\subsection{Future Work}

Giovannini et. al have contributed an Agda mechanization of most of the concrete
model construction described in their paper. % TODO citation?
%
It would be interesting to mechanize the abstract model we have formulated in
this work and instantiate it to obtain the concrete model they constructed as
well as models for languages with other effects as described in this work.


\bibliographystyle{plain}
\bibliography{../../paper-new/references}

\appendix

\section{Call-by-push-value}

\subsection{Morphisms of CBPV Models}

There are two relevant notions of \emph{morphism} of CBPV models:
\emph{strict} and \emph{lax}.
% morphisms of CBPV models
Given call-by-push-value models
$\mathcal M_1 = (\mathcal V_1, \mathcal E_1, \arr_1, U_1, F_1)$ and
$\mathcal M_2 = (\mathcal V_2, \mathcal E_2, \arr_2, U_2, F_2)$,
A \emph{strict} morphism $G$ from $M_1$ to $M_2$ is given by a pair of functors
$G_{\mathcal{V}}: V_1 \to V_2$ and $G_{\mathcal{E}} : E_1 \to E_2$
that strictly presere the constructors:
\begin{enumerate}
  \item $G_{\mathcal{E}} \circ F_1 = F_2 \circ G_{\mathcal{V}}$
  \item $G_{\mathcal{V}} \circ U_1 = U_2 \circ G_{\mathcal{E}}$
  \item $G_{\mathcal{E}}(A \arr_1 B) = G_{\mathcal{V}}(A) \arr_2 G_{\mathcal{E}}(B)$
  \item $G_{\mathcal{V}}(A_1 \times_1 A_2) = G_{\mathcal{V}}(A_1) \times_2 G_{\mathcal{V}}(A_2)$
  \item $G_{\mathcal{V}}(1_1) = 1_2$
\end{enumerate}
As well as strictly preserving the corresponding universal morphisms
and coherence isomorphisms.

A lax morphism instead preserves these only up to transformation
\begin{enumerate}
  \item $G_{\mathcal{E}} \circ F_1 \Rightarrow F_2 \circ G_{\mathcal{V}}$
  \item $G_{\mathcal{V}} \circ U_1 \Rightarrow U_2 \circ G_{\mathcal{E}}$
  \item $G_{\mathcal{E}}(A \arr_1 B) \Rightarrow G_{\mathcal{V}}(A) \arr_2 G_{\mathcal{E}}(B)$
  \item $G_{\mathcal{V}}(A_1 \times_1 A_2) \Rightarrow G_{\mathcal{V}}(A_1) \times_2 G_{\mathcal{V}}(A_2)$
  \item $G_{\mathcal{V}}(1_1) \Rightarrow 1_2$
\end{enumerate}
Additionally a lax morphism should have a relationship between these
transformations and the universal morphisms, but we will only consider
lax morphisms of thin categories, where such conditions hold
trivially.

\subsection{Kleisli Actions of CBPV Type Constructors}

In CBPV models, all the type constructors are interpreted as functors:
\begin{enumerate}
\item $\to : \op\calV \times \calE \to \calE$
\item $\times : \calV \times \calV \to \calV$
\item $F : \calV \to \calE$
\item $U : \calE \to \calV$
\end{enumerate}
That is, they all have functorial actions on \emph{pure} morphisms of
value types and \emph{linear} morphisms of computation types.
%
We use these functorial actions extensively in the construction of
casts and their corresponding perturbations. But when defining
downcasts of value types and upcasts of computation types, we
additionally need a second functorial action of these categories:
functoriality in \emph{impure} morphisms of value types and
\emph{non-linear} morphisms of computation types. These notions of
morphism are given by the \emph{Kleisli} categories $\calVk$ and
$\calEk$ which have value types and computation types as objects but
morphisms are defined as
\[ \calVk(A,A') = \calE(F A, FA')\]
\[ \calEk(B,B') = \calV(U B, U B')\]
with composition given by composition in $\calE/\calV$.  That is we
need to define a second functorial action, that agrees with the above
on objects for these Kleisli categories:
\begin{enumerate}
\item $\tok : \op\calVk \square \calEk \to \calEk$
\item $\timesk : \calVk \square \calVk \to \calVk$
\item $\Fk : \calVk \to \calEk$
\item $\Uk : \calEk \to \calVk$
\end{enumerate}
Note that rather than the product of categories we use the ``funny
tensor product'' $\square$. This is because the action on
impure/non-linear morphisms for $\tok/\timesk$ do not satisfy ``joint
functoriality'' but instead only ``separate functoriality'', meaning
we give rather than an action on morphisms in both categories
simultaneously instead an action on each argument categories morphisms
with the object in the other category fixed. The existence of these
functorial actions for $\tok$ and $\timesk$ is reliant on the
\emph{strength} of the adjunction. We describe them using the internal
language of CBPV in order to more easily verify their
existence/functoriality:
\begin{enumerate}
\item For $\tok$ we define for $\phi : \calE(F A,F A')$ and $B \in \calE$ the morphism $\phi \tok B : \calV(U(A' \to B),U(A\to B))$ as
  \[ t:U(A'\to B) \vdash \phi \tok B = \{ \lambda x. x' \leftarrow \phi\,[\ret x]; ! t x'\} : U(A \to B) \]
  and for $A \in \calV$ and $f : \calV(UB,UB')$ we define $A \tok f : \calV(U(A \to B),U(A\to B'))$ as
  \[ t : U(A \to B) \vdash A \tok f = \{ \lambda x. !f[\{ ! t x \}]\} \]
\item For $\timesk$ we define for $\phi : \calE(F A_1,FA_2)$ and $A' \in \calV$ the morphism $\phi \timesk A_2$ as
  \[ \bullet : F(A_1\times A_2) \vdash \phi \timesk A_2 = (x_1,x_2) \leftarrow \bullet; x_1' \leftarrow \phi[\ret x_1]; \ret (x_1',x_2) : F(A_1'\times A_2)\]
  and $A_1 \timesk \phi$ is defined symmetrically.
\item For $\Uk$ we need to define for $f : \calV(UB,UB')$ a morphism $\Uk f : \calE(FUB,FUB')$. This is simply given by the functorial action of $F$: $\Uk f = F(f)$
\item Similarly $\Fk \phi = U\phi$
\end{enumerate}

Functoriality in each argument is easily established, meaning for
example for the function type is functorial in each argument:
\begin{enumerate}
\item $(\phi \circ \phi') \tok B = (\phi' \tok B) \circ (\phi \tok B)$
\item $\id \tok B = \id$
\item $A \tok (f \circ f') = (A \tok f) \circ (A \tok f)$
\item $A \tok \id = \id$
\end{enumerate}

Finally, note that all of these constructions lift to squares in a
double CBPV model since the squares themselves form a CBPV model and
the projection functions preserve CBPV structure. For instance, given a square
$\alpha : \phi \ltdyn_{F c_o}^{F c_i} \phi'$ and a horizontal morphism $d : B \rel B'$ of appropriate type, we get a square
\[ \alpha \tok d : \phi \tok B \ltdyn_{U(c_o \to d)}^{U(c_i \to d)} \phi' \tok B' \]

\end{document}